\documentclass[portuguese,12pt,oneside,a4paper]{book}

% ---------------------------------------------------------------------------- %
% Pacotes
\usepackage[T1]{fontenc}
\usepackage[brazil]{babel}           % define a linguagem padrão
\usepackage[utf8x]{inputenc}
\usepackage[scaled]{helvet}          % Arial (via Helvetica, equivalente Type1) — Quadro 2.3.3 do manual IFPB-CG
\renewcommand{\familydefault}{\sfdefault}
\usepackage[pdftex]{graphicx}           % usamos arquivos pdf/png como figuras
\usepackage{setspace}                   % espaçamento flexível
\usepackage{makeidx}                    % índice remissivo
\usepackage[nottoc,notlof,notlot]{tocbibind} % NÃO acrescenta lista de figuras/tabelas/sumário ao próprio Sumário
\usepackage{courier}                    % usa o Adobe Courier no lugar de Computer Modern Typewriter
\usepackage{type1cm}                    % fontes realmente escaláveis
\usepackage{listings}                   % para formatar código-fonte (ex. em Java)
\usepackage{titletoc}
\usepackage[fixlanguage]{babelbib}
\usepackage[font=small,format=plain,labelfont=bf,up,textfont=it,up]{caption}
\usepackage[usenames,svgnames,dvipsnames]{xcolor}
\usepackage[a4paper,top=2.54cm,bottom=2.0cm,left=2.0cm,right=2.54cm]{geometry} % margens
\usepackage[pdftex,plainpages=false,pdfpagelabels,pagebackref,colorlinks=true,citecolor=black,linkcolor=black,urlcolor=black,filecolor=black,bookmarksopen=true]{hyperref} % links em preto
\usepackage[all]{hypcap}                % soluciona o problema com o hyperref e capitulos
\usepackage{indentfirst}                % recuo também no primeiro parágrafo após título (norma ABNT)

\usepackage{color} % colorir texto
\usepackage{xcolor}
\usepackage{amsmath} %expressoes matematicas
\usepackage{pdflscape} %permite mudar orientacao da pagina para paisagem E MUDA NO PDF
\usepackage{longtable} %permite tabelas que ocupem mais de uma página
\usepackage{multicol}
\usepackage{enumitem} %para pode iniciar enumerate em 0
\usepackage{pdfpages}
\usepackage{tikz}                       % diagramas do capítulo de implementação
\usetikzlibrary{arrows.meta,positioning,shapes.geometric,fit,calc,backgrounds}


% ---------------------------------------------------------------------------- %
% Formatação de seções conforme Quadro 2 do Manual de Elaboração de Trabalhos
% Acadêmicos do IFPB Campus Campina Grande (NBR 6024): primária em negrito e
% maiúsculo; secundária em maiúsculo sem negrito; terciária com negrito e
% inicial maiúscula; quaternária sem negrito e inicial maiúscula. O quadro não
% pede tamanho de fonte maior para nenhum nível — capítulo, seção e texto
% ficam todos no mesmo tamanho (12pt), diferenciados só por negrito/caixa alta.
% O espaçamento entre título e texto — e entre o fim de uma seção e o título
% seguinte — é de apenas uma linha em branco no mesmo espaçamento do texto
% (1,5), não um vão maior.
\usepackage{titlesec}
\titleformat{\chapter}[hang]{\normalfont\bfseries}{\thechapter}{1em}{\MakeUppercase}
\titleformat{name=\chapter,numberless}[block]{\normalfont\bfseries\centering}{}{0pt}{}
\titleformat{\section}{\normalfont}{\thesection}{1em}{\MakeUppercase}
\titleformat{\subsection}{\normalfont\bfseries}{\thesubsection}{1em}{}
\titleformat{\subsubsection}{\normalfont}{\thesubsubsection}{1em}{}
\titlespacing*{\chapter}{0pt}{0pt}{1\baselineskip}
\titlespacing*{\section}{0pt}{1\baselineskip}{1\baselineskip}
\titlespacing*{\subsection}{0pt}{1\baselineskip}{1\baselineskip}
\titlespacing*{\subsubsection}{0pt}{1\baselineskip}{1\baselineskip}

% ---------------------------------------------------------------------------- %
\graphicspath{{./figuras/}}             % caminho das figuras (recomendável)
\frenchspacing                          % arruma o espaço: id est (i.e.) e exempli gratia (e.g.)
\urlstyle{same}                         % URL com o mesmo estilo do texto e não mono-spaced
\makeindex                              % para o índice remissivo
\raggedbottom                           % para não permitir espaços extra no texto

% ---------------------------------------------------------------------------- %
% Opções de listing usados para o código fonte
% Ref: http://en.wikibooks.org/wiki/LaTeX/Packages/Listings
\lstset{ %
language=Java,                  % choose the language of the code
basicstyle=\footnotesize,       % the size of the fonts that are used for the code
numbers=left,                   % where to put the line-numbers
numberstyle=\footnotesize,      % the size of the fonts that are used for the line-numbers
stepnumber=1,                   % the step between two line-numbers. If it's 1 each line will be numbered
numbersep=5pt,                  % how far the line-numbers are from the code
showspaces=false,               % show spaces adding particular underscores
showstringspaces=false,         % underline spaces within strings
showtabs=false,                 % show tabs within strings adding particular underscores
frame=single,	                % adds a frame around the code
framerule=0.6pt,
tabsize=2,	                    % sets default tabsize to 2 spaces
captionpos=b,                   % sets the caption-position to bottom
breaklines=true,                % sets automatic line breaking
breakatwhitespace=false,        % sets if automatic breaks should only happen at whitespace
escapeinside={\%*}{*)},         % if you want to add a comment within your code
backgroundcolor=\color[rgb]{1.0,1.0,1.0}, % choose the background color.
rulecolor=\color[rgb]{0.8,0.8,0.8},
extendedchars=true,
xleftmargin=10pt,
xrightmargin=10pt,
framexleftmargin=10pt,
framexrightmargin=10pt
}

\lstdefinelanguage{blah} % minha linguagem se chama blah
{
% Aqui vão os atributos dela. Para ver mais atributos, Google!
tabsize=2, % o tab possui 2 espaços
frame=shadowbox,% o tipo de box que usaremos
%rulesepcolor=\color{lightgrey}, %/ a cor pra fazer o efeito legal da box
captionpos=b,% o caption do código vai para baixo (botton) o default eh top
extendedchars=false, % eu usei isso para conseguir colocar expressoes matematicas utilizando o scapeinside tbm
escapeinside='', % usando aspas, pode-se usar qualquer comando dentro deles
morekeywords = {if, else, while, do, for, and, or, action, initialize} % palavras reservadas que ficaram em negrito
}


% ABNT %%%%%%%%%%%%%%%%%%%%%%%%%%%%%%%%%%%%%%%%%%%%%%%%%%%%%%%%%%%%%%%%%%%%%%%%%%%%%
\usepackage[alf,% Para usar o estilo de citacao alfabetico (seção 9.2 em 6023/2000
abnt-etal-text=emph,% o et al. vai aparecer em italico
abnt-and-type=e,% separador de autores com ´e´
bibjustif,% bibliografia justificada no final
abnt-etal-cite=3]{abntex2cite}
% No caso, a 6023 diz que se deve usar ‘et al.’ para mais de três autores, ou seja, a partir de quatro autores. Então, usarei abnt-etal-cite=3
\usepackage{cite}%%% NOVO

% ------------ colocar o [fulanclearo 2002] antes da descricao de referencia bibliografica
\makeatletter
\renewcommand\@biblabel[1]{[#1]~}% colocar citação na frente da lista de bibliografia
\renewcommand\citepunct{;~}%%% NOVO
%\renewcommand{\@biblabel}[1]{[\textbf{#1}]~}
\makeatother
\citebrackets[] % fazer com que as citacoes dentro do texto virem colchetes
% -------------------------------------------------------
%%%%%%%%%%%%%%%%%%%%%%%%%%%%%%%%%%%%%%%%%%%%%%%%%%%%%%%%%%%%%%%%%%%%%%%%%%%%%%%%%%%%


% Entrada de apêndice no Sumário: "Apêndice A --- Título" (sem largura fixa de
% número). Cada arquivo de apêndice liga \apendiceNumberline antes do \chapter e
% restaura \numberlineorig depois, para que as seções sigam o formato padrão.
\let\numberlineorig\numberline
\newcommand{\apendiceNumberline}[1]{#1\ ---\ }

% ---------------------------------------------------------------------------- %
% Corpo do texto
\begin{document}
\frontmatter
\setcounter{tocdepth}{2}

% Paginação (seção 2.3.4 do manual IFPB-CG): as páginas são contadas a partir
% da folha de rosto, mas só começam a ser graficamente numeradas na primeira
% página da Introdução. Por isso os elementos pré-textuais ficam sem número
% visível (pagestyle vazio); a numeração aparece só a partir do \mainmatter.
\makeatletter
\let\ps@plainorig\ps@plain
\renewcommand{\ps@plain}{\ps@empty}
\makeatother
\pagestyle{empty}

\onehalfspacing  % espaçamento

% ---------------------------------------------------------------------------- %
% Capa

\thispagestyle{empty}
\begin{center}
    \includegraphics[scale=0.4]{ifpb} \\
    INSTITUTO FEDERAL DE EDUCAÇÃO, CIÊNCIA E TECNOLOGIA DA PARAÍBA\\
    CAMPUS CAMPINA GRANDE\\
    COORDENAÇÃO DO CURSO SUPERIOR BACHARELADO EM ENGENHARIA DE COMPUTAÇÃO\\

    \vspace*{4cm}

    \textbf{GABRIEL DOMINGOS DE MEDEIROS}

    \vspace*{3.5cm}
    \textbf{IMPLEMENTAÇÃO E VALIDAÇÃO DE UMA MÁQUINA VIRTUAL DE LÓGICA LADDER E AGNÓSTICA DE HARDWARE BASEADA EM MODEL-BASED DESIGN}\\

    \vfill

    Campina Grande - PB\\
    2026

\end{center}

\newpage

\thispagestyle{empty}
\begin{center}

    \textbf{GABRIEL DOMINGOS DE MEDEIROS}

    \vspace*{2cm}
    \textbf{IMPLEMENTAÇÃO E VALIDAÇÃO DE UMA MÁQUINA VIRTUAL DE LÓGICA LADDER E AGNÓSTICA DE HARDWARE BASEADA EM MODEL-BASED DESIGN}\\

    \vskip 1.5cm
\end{center}



    \begin{flushright}
		\begin{minipage}[left]{0.5\linewidth}

			Monografia apresentada à Coordenação do Curso de Engenharia de Computação do IFPB - Campus Campina Grande, como requisito parcial para conclusão do curso de Bacharelado em Engenharia de Computação.\\

		\end{minipage}

   	\end{flushright}


\begin{center}

    \vskip 1cm
    Instituto Federal da Paraíba

    \vskip 2cm
    Orientador: Fagner de Araújo Pereira\\

    \vfill

    Campina Grande - PB\\
    2026
\end{center}

\newpage

% ---------------------------------------------------------------------------- %
% FICHA CATALOGRÁFICA (VERSÃO FINAL/CORRIGDA)
% ---------------------------------------------------------------------------- %

%\includepdf{ficha-catalografica.pdf} %Incluir a ficha fornecida pela biblioteca do IFPB


% ---------------------------------------------------------------------------- %
% Página de rosto (só para a versão final)
% ---------------------------------------------------------------------------- %

\newpage
\thispagestyle{empty}
    \begin{center}


	    \textbf{GABRIEL DOMINGOS DE MEDEIROS}

        \vspace*{2cm}

    	\textbf{IMPLEMENTAÇÃO E VALIDAÇÃO DE UMA MÁQUINA VIRTUAL DE LÓGICA LADDER E AGNÓSTICA DE HARDWARE BASEADA EM MODEL-BASED DESIGN}\\

    	\vskip 2cm
    \end{center}

    \begin{center}
		\begin{minipage}{0.75\linewidth}

			Monografia apresentada ao Curso de Bacharelado em Engenharia de Computação, do Instituto Federal de Educação, Ciência e Tecnologia da Paraíba - Campus Campina Grande, em cumprimento às exigências parciais para a obtenção do título de Bacharel em Engenharia da Computação.\\

			\vskip 0.5cm

			Trabalho aprovado. Campina Grande - PB, \rule{1cm}{0.4pt} de \rule{3cm}{0.4pt} de 2026.

		\end{minipage}
    \end{center}

    \begin{center}

    	\vskip 1cm

		\rule{10cm}{0.1mm}\\
		FAGNER DE ARAÚJO PEREIRA\\
		Orientador

    	\vskip 2cm

		\rule{10cm}{0.1mm}\\
		XXXXXXXXXXXXXXXXXXXXXXXXXXXXXXX\\
		Membro da Banca

    	\vskip 2cm

		\rule{10cm}{0.1mm}\\
		XXXXXXXXXXXXXXXXXXXXXXXXXXXXXXXX\\
		Membro da Banca

		\vfill

		Campina Grande - PB\\
		2026

    \end{center}
\pagebreak

% % ---------------------------------------------------------------------------- %
% % Agradecimentos

\chapter*{Agradecimentos}


	XXXXXXXXXXXXXXXXXXXXXX


\pagebreak

\newpage
\thispagestyle{empty}


        \vspace*{20.0cm}

        \begin{flushright}
			\textit{Dedico este trabalho à minha família, pelo apoio incondicional, e a todos que, de alguma forma, contribuíram para que este momento fosse possível.}
        \end{flushright}



\pagebreak

\newpage
\thispagestyle{empty}


        \vspace*{20.0cm}

        \begin{flushright}

             Alguma frase/pensamento (OPCIONAL)\\
             Autor\\


        \end{flushright}



\pagebreak


% ---------------------------------------------------------------------------- %
% Resumo

\chapter*{Resumo}
	Este trabalho apresenta a implementação e a validação de uma máquina virtual de lógica \textit{Ladder} independente de \textit{hardware}, concebida segundo a metodologia de \textit{Model-Based Design} para executar programas de controle em microcontroladores. A solução compreende um conjunto de 76 instruções com palavra-base de 32 \textit{bits}, um motor de execução com memória alocada estaticamente e acessos validados, uma camada de abstração que isola o motor do microcontrolador e um protocolo serial, com enquadramento COBS e verificação CRC16, para a carga do \textit{bytecode} e a depuração. A plataforma-alvo é o microcontrolador STM32F767ZI, e a máquina de modos do bloco de controle PID é especificada em \textit{Simulink}/\textit{Stateflow}. A validação seguiu um catálogo de 25 casos de teste, que cobre lógica de contatos, temporização, contagem, aritmética, entradas e saídas analógicas, blocos de controle, robustez do protocolo e desempenho, com resultados esperados obtidos de tabelas-verdade, de cálculos independentes e de um modelo de referência executável. Dos 220 passos de verificação funcional, 219 foram conformes, e os 52 pacotes malformados aplicados ao protocolo foram rejeitados com o código de rejeição esperado. Os ensaios revelaram quatro defeitos de implementação, posteriormente corrigidos. Na configuração de depuração do compilador, o tempo de varredura foi de cerca de 37~$\mu$s por \textit{rung} de cinco instruções, e o \textit{firmware} ocupou 76,67~KB de memória \textit{Flash} e 15,6~KB de RAM. A validação cruzada contra os modelos \textit{Simulink}/\textit{Stateflow} e a repetição dos ensaios em malha fechada na plataforma-alvo permanecem como trabalhos futuros.\\

\noindent \textbf{Palavras-chave:} Máquina virtual. Lógica \textit{Ladder}. \textit{Model-Based Design}. Controlador lógico programável. STM32.


\pagebreak


% ---------------------------------------------------------------------------- %

 \chapter*{Abstract}
This work presents the implementation and validation of a hardware-agnostic Ladder logic virtual machine, designed under the Model-Based Design methodology to run control programs on microcontrollers. The solution comprises a 76-instruction set with a 32-bit base word, an execution engine with statically allocated memory and validated accesses, a hardware abstraction layer that isolates the engine from the microcontroller, and a serial protocol, with COBS framing and CRC16 checking, for bytecode download and debugging. The target platform is the STM32F767ZI microcontroller, and the mode logic of the PID control block is specified in Simulink/Stateflow. Validation followed a catalogue of 25 test cases covering contact logic, timing, counting, arithmetic, analog inputs and outputs, control blocks, protocol robustness and performance, with expected results taken from truth tables, independent calculations and an executable reference model. Of the 220 functional verification steps, 219 were conforming, and all 52 malformed packets applied to the protocol were rejected with the expected rejection code. The tests revealed four implementation defects, which were subsequently fixed. In the compiler debug configuration, the scan time was about 37~$\mu$s per five-instruction rung, and the firmware occupied 76.67~KB of Flash memory and 15.6~KB of RAM. Cross-validation against the Simulink/Stateflow models and the repetition of the closed-loop tests on the target platform remain as future work.\\

 \noindent \textbf{Keywords:} Virtual machine. Ladder logic. Model-Based Design. Programmable logic controller. STM32.



% ---------------------------------------------------------------------------- %
% Listas de figuras e tabelas criadas automaticamente
\listoffigures
\listofquadros
\listoftables


% ---------------------------------------------------------------------------- %
\chapter*{Lista de Abreviaturas}
\begin{tabular}{ll}
	ADC		& \textit{Analog-to-Digital Converter} \\
	ARM 	& \textit{Advanced RISC Machine} \\
	CLP		& Controlador Lógico Programável \\
	COBS	& \textit{Consistent Overhead Byte Stuffing} \\
	CRC		& \textit{Cyclic Redundancy Check} \\
	DMA		& \textit{Direct Memory Access} \\
	FBD		& \textit{Function Block Diagram} \\
	FPU		& \textit{Floating-Point Unit} \\
	GPIO	& \textit{General Purpose Input/Output} \\
	HAL		& \textit{Hardware Abstraction Layer} \\
	HIL		& \textit{Hardware-in-the-Loop} \\
	IEC		& \textit{International Electrotechnical Commission} \\
	IL		& \textit{Instruction List} \\
	IoT		& \textit{Internet of Things} \\
	ISA		& \textit{Instruction Set Architecture} \\
	LD		& \textit{Ladder Diagram} \\
	MBD		& \textit{Model-Based Design} \\
	PID		& Proporcional-Integral-Derivativo \\
	POSIX	& \textit{Portable Operating System Interface} \\
	PWM		& \textit{Pulse-Width Modulation} \\
	RAM		& \textit{Random Access Memory} \\
	SFC		& \textit{Sequential Function Chart} \\
	ST		& \textit{Structured Text} \\
	USART	& \textit{Universal Synchronous/Asynchronous Receiver/Transmitter} \\
	VM		& \textit{Virtual Machine} (Máquina Virtual) \\
	vPLC	& \textit{Virtual Programmable Logic Controller} \\

\end{tabular}



% ---------------------------------------------------------------------------- %
% Sumário
\tableofcontents    % imprime o sumário


% ---------------------------------------------------------------------------- %
% Capítulos do trabalho
\mainmatter

% Paginação volta a ficar visível (numeração arábica) a partir daqui, conforme
% seção 2.3.4 do manual IFPB-CG.
\makeatletter
\let\ps@plain\ps@plainorig
\makeatother
\pagestyle{plain}

\onehalfspacing            % espaçamento um e meio

% !TeX root = ../main.tex
%% ------------------------------------------------------------------------- %%
\chapter{Introdução}
\label{cap:introducao}

%% ------------------------------------------------------------------------- %%
\section{Contextualização e Descrição do Problema}
\label{sec:contextualizacao}

O mercado global de Controladores Lógicos Programáveis (CLPs) tem apresentado crescimento expressivo nos últimos anos, impulsionado pelos investimentos crescentes em automação industrial e pela modernização de sistemas de controle legados: estima-se que o setor movimenta atualmente cerca de US\$ 13,33 bilhões, com projeção de expansão contínua até 2031 \cite{mordor:2026}.
Esse crescimento reflete a consolidação da automação industrial como pilar central da Indústria 4.0, na qual controladores programáveis assumem papel crítico na interconexão de máquinas, sensores e processos produtivos.
Nesse cenário, contudo, a expansão do setor convive com uma forte fragmentação tecnológica: a predominância de soluções proprietárias e de arquiteturas de \textit{hardware} específicas por fabricante ainda limita a portabilidade de algoritmos de controle e favorece o aprisionamento tecnológico dos usuários finais.
É nesse contexto que se insere a proposta deste trabalho: uma máquina virtual (\textit{VM}) embarcada, capaz de interpretar instruções de lógica \textit{Ladder} --- baseada na norma IEC 61131-3 \cite{iec:2025} --- de maneira independente da plataforma de \textit{hardware} subjacente.

A lógica \textit{Ladder} (\textit{Ladder Diagram}) é uma das cinco linguagens de programação padronizadas pela IEC 61131-3 para controladores industriais, caracterizada pela representação gráfica de contatos e bobinas que remete aos antigos diagramas elétricos a relés, o que a tornou historicamente a linguagem mais difundida na automação industrial.
Uma máquina virtual, por sua vez, corresponde a uma camada de \textit{software} capaz de interpretar um conjunto de instruções (\textit{bytecode}) de forma independente do \textit{hardware} subjacente --- de modo análogo ao papel desempenhado por máquinas virtuais de propósito geral em outros domínios da computação.
Reunir esses dois conceitos exige a definição de uma Arquitetura de Conjunto de Instruções (\textit{ISA} --- \textit{Instruction Set Architecture}) customizada e de tamanho fixo (32 \textit{bits}), projetada para otimizar o processo de busca e decodificação (\textit{fetch-decode}) e garantir o determinismo do tempo de varredura; a pesquisa foca especificamente na migração dessa arquitetura para microcontroladores da família \textit{STM32} (núcleo \textit{ARM Cortex-M7}), explorando técnicas de abstração de \textit{hardware} (\textit{HAL}) e metodologias de \textit{Model-Based Design}.
Propõe-se ainda a segregação estrutural entre o motor de execução (\textit{firmware} estático) e a lógica de aplicação (\textit{bytecode} tratado como dado), estabelecendo um ambiente de execução isolado (\textit{sandbox}) que previne falhas críticas de processamento e viabiliza a gestão segura do ciclo de vida da aplicação industrial.

Apesar do crescimento do setor, a rigidez das arquiteturas convencionais de CLPs impõe barreiras significativas à inovação, dificultando a portabilidade de algoritmos de controle para plataformas de \textit{hardware} mais modernas e acessíveis sem a necessidade de reengenharia completa do \textit{firmware}.
Além disso, a validação de lógicas de controle complexas diretamente no \textit{hardware} final apresenta riscos operacionais e custos elevados, carecendo de metodologias que integrem eficientemente a simulação computacional com a execução física.

No âmbito das soluções de código aberto atualmente disponíveis, como o ecossistema \textit{OpenPLC} \cite{alves:2018}, a abordagem predominante baseia-se na compilação de diagramas \textit{Ladder} para código C nativo, seguida de compilação estática.
Embora eficiente em termos de tempo de execução, essa metodologia funde a lógica de controle aos \textit{drivers} de \textit{hardware} em um binário monolítico, de modo que a adaptação para novas plataformas ou pinagens exige modificações diretas no código-fonte dos \textit{drivers} em baixo nível, limitando drasticamente a portabilidade.
Adicionalmente, a execução do código traduzido diretamente no processador elimina barreiras de proteção inerentes a sistemas operacionais: um erro na modelagem da lógica do usuário pode resultar em falhas críticas do sistema, comprometendo a estabilidade do equipamento industrial.

Por outro lado, abordagens alternativas baseadas em interpretadores de linguagens de \textit{script} genéricas, como o \textit{MicroPython} \cite{micropython:2025}, falham em atender aos requisitos de tempo real estrito da automação industrial, devido à gestão não determinística de memória (\textit{Garbage Collection}), que introduz variações inaceitáveis no tempo de varredura.
Diante do exposto, o problema técnico enfrentado por este trabalho consiste na ausência de uma camada de virtualização eficiente e aberta que permita a execução segura e determinística de lógicas de controle padronizadas em microcontroladores de propósito geral, garantindo a fidelidade comportamental entre o modelo projetado em ambiente de simulação e o sistema embarcado final.


%% ------------------------------------------------------------------------- %%
\section{Objetivos}
\label{sec:objetivo}

\subsection{Objetivo Geral}
\label{sec:objgeral}

Diante do problema apresentado, este trabalho tem como objetivo geral desenvolver e validar uma máquina virtual embarcada capaz de interpretar \textit{bytecode} de lógica \textit{Ladder} com fidelidade semântica à norma IEC 61131-3, executando de forma determinística e agnóstica de \textit{hardware} em um microcontrolador da família \textit{STM32}.

\subsection{Objetivos Específicos}
\label{sec:objespecificos}

\begin{itemize}
	\item Especificar e documentar formalmente uma Arquitetura de Conjunto de Instruções (\textit{ISA}) de 32 \textit{bits}, compatível com as primitivas operacionais do subconjunto \textit{Ladder Diagram} da IEC 61131-3;
	\item Modelar lógicas de controle de referência em ambiente \textit{Simulink}/\textit{Stateflow}, segundo a metodologia de \textit{Model-Based Design};
	\item Implementar, em linguagem C, o motor de execução (\textit{VM}) responsável pelo ciclo de busca, decodificação e execução do \textit{bytecode};
	\item Desenvolver uma Camada de Abstração de \textit{Hardware} (\textit{HAL}) para o microcontrolador \textit{STM32F7} (núcleo \textit{ARM Cortex-M7});
	\item Validar cruzadamente o comportamento da \textit{VM} embarcada frente ao modelo de referência simulado, comparando as tabelas de imagem de saída para um mesmo conjunto de estímulos de entrada;
	\item Analisar quantitativamente o desempenho da solução, mensurando o tempo de varredura (\textit{scan cycle}) e o consumo de memória \textit{RAM} e \textit{Flash}.
\end{itemize}


%% ------------------------------------------------------------------------- %%
\section{Justificativa}
\label{sec:justificativa}

A relevância deste trabalho decorre da lacuna identificada entre as soluções de código aberto disponíveis e os requisitos de portabilidade e determinismo da automação industrial.
Ao segregar o motor de execução da lógica de aplicação, a arquitetura proposta elimina a necessidade de recompilação do \textit{firmware} a cada alteração da lógica \textit{Ladder}, reduzindo o tempo de comissionamento e o risco de falhas decorrentes de atualizações completas de memória.
O caráter de ambiente isolado (\textit{sandbox}) confere ainda uma camada de segurança inexistente nas soluções que traduzem a lógica do usuário diretamente para código nativo, contendo eventuais erros estruturais antes que comprometam a estabilidade do sistema embarcado.

Do ponto de vista social e econômico, a proposta de uma solução aberta e portável para múltiplas plataformas de \textit{hardware} contribui para reduzir o aprisionamento tecnológico a fabricantes específicos, democratizando o acesso a tecnologias de automação hoje concentradas em soluções proprietárias de alto custo.


%% ------------------------------------------------------------------------- %%
\section{Estrutura do Documento}
\label{sec:organizacao}

A metodologia adotada caracteriza-se como pesquisa aplicada de natureza experimental e abordagem predominantemente quantitativa, estruturada em etapas sequenciais com entregas incrementais que partem da especificação formal da arquitetura até a validação funcional integrada ao \textit{hardware}; o detalhamento de cada etapa --- especificação da \textit{ISA}, modelagem de referência, implementação do motor de execução, desenvolvimento da \textit{HAL}, integração com \textit{hardware}, validação cruzada e análise de desempenho --- será apresentado nos capítulos de desenvolvimento deste trabalho.

O restante deste trabalho está organizado da seguinte forma: o Capítulo~\ref{cap:fundamentacao} apresenta a fundamentação teórica, contemplando a norma IEC 61131-3, os trabalhos relacionados e o estado da prática em soluções para execução de lógicas de controle industrial, além dos conceitos de arquitetura que sustentam a solução proposta.
O Capítulo~\ref{cap:implementacao} descreve a implementação --- arquitetura, plataforma-alvo, máquina virtual e conjunto de instruções, modelo de memória, camada de abstração de \textit{hardware}, protocolo de comunicação e modelagem de referência --- com a justificativa de cada escolha de projeto, e apresenta a estratégia de validação, o catálogo de testes e os resultados dos ensaios sobre a plataforma-alvo.
Seguem-se as considerações finais e as sugestões para trabalhos futuros, no Capítulo~\ref{cap:finais}.
A especificação completa do conjunto de instruções, os ensaios do bloco de controle PID e o registro do ensaio de robustez do protocolo constam dos apêndices.

% !TeX root = ../main.tex
%% ------------------------------------------------------------------------- %%
\chapter{Fundamentação Teórica}
\label{cap:fundamentacao}

%% ------------------------------------------------------------------------- %%
\section{A Norma IEC 61131-3}
\label{sec:iec61131}

A norma IEC 61131-3, cuja quarta edição foi publicada em maio de 2025 pela Comissão Eletrotécnica Internacional (IEC), constitui a norma internacional que define a arquitetura de \textit{software} e as linguagens de programação para controladores programáveis \cite{iec:2025}.
A norma especifica a sintaxe e a semântica de um conjunto unificado de linguagens, incluindo \textit{Structured Text} (\textit{ST}), \textit{Ladder Diagram} (\textit{LD}), \textit{Function Block Diagram} (\textit{FBD}) e \textit{Sequential Function Chart} (\textit{SFC}).
A edição mais recente removeu a linguagem \textit{Instruction List} (\textit{IL}), anteriormente marcada como obsoleta, e incorporou suporte a \textit{strings UTF-8}.
A importância da norma para o presente trabalho é fundamental, pois a \textit{ISA} customizada da máquina virtual proposta é projetada para representar fielmente as primitivas operacionais definidas no subconjunto \textit{LD} da IEC 61131-3, garantindo compatibilidade semântica com o padrão industrial.


%% ------------------------------------------------------------------------- %%
\section{Trabalhos Relacionados}
\label{sec:trabalhos-relacionados}

No domínio de execução portável de lógicas de controle industrial, destacam-se três eixos temáticos principais na literatura: a implementação do modelo de \textit{software} da norma IEC 61131-3 em sistemas operacionais de tempo real, o desenvolvimento de controladores lógicos de código aberto e as metodologias de \textit{Model-Based Design} aplicadas à geração automática de código para CLPs.

\subsection{Implementação da IEC 61131-3 sobre Sistemas POSIX de Tempo Real}
\label{sub:plaza-medrano}

Plaza e Medrano \cite{plaza:2006} apresentaram uma proposta de implementação do modelo de \textit{software} definido pela norma IEC 61131-3 utilizando sistemas operacionais compatíveis com \textit{POSIX} (\textit{Pthreads} e temporizadores de alta resolução), especificamente o \textit{RTLinux}.
Os autores discutem os aspectos críticos da norma que necessitam de esclarecimento antes de sua implementação prática, como concorrência, escopo de variáveis e operação cíclica, demonstrando que o modelo IEC 61131-3 pode ser mapeado para primitivas \textit{POSIX}, gerando código C a partir de descrições textuais em \textit{Instruction List} (\textit{IL}).
A contribuição central do trabalho reside na comprovação de que a camada de tradução de IEC para C no padrão \textit{POSIX} é viável e que o desempenho resultante é comparável ao de produtos comerciais.
Por outro lado, a abordagem permanece vinculada a plataformas com sistema operacional \textit{Linux} em tempo real, não abordando a execução em microcontroladores \textit{bare-metal} de recursos limitados, cenário de aplicação do presente trabalho.

\subsection{Controladores Lógicos de Código Aberto}
\label{sub:openplc-alves}

Alves \cite{alves:2018} apresentou o desenvolvimento do \textit{OpenPLC}, o primeiro CLP de código aberto totalmente compatível com a norma IEC 61131-3.
O trabalho detalha a arquitetura interna do sistema, composta pelo compilador \textit{MatIEC} --- responsável por traduzir código IEC 61131-3 para C++ ---, servidor \textit{web}, camadas de rede, tabelas de imagem de E/S e uma camada de \textit{hardware}.
Essa obra contribui significativamente para a identificação de limitações arquiteturais do \textit{OpenPLC}: a compilação estática via \textit{MatIEC} funde a lógica do usuário ao binário do sistema, eliminando a possibilidade de atualização dinâmica da lógica de controle e introduzindo riscos de falhas catastróficas caso a lógica do usuário contenha erros estruturais.
Essas limitações motivam diretamente a arquitetura de segregação entre motor de execução e \textit{bytecode} proposta neste trabalho.

\subsection{Geração e Validação Automática de Código via Model-Based Design}
\label{sub:duggan}

Duggan et al. \cite{duggan:2022} apresentaram uma metodologia para projeto e validação de equipamentos em ambientes de manufatura regulados, utilizando o \textit{Simulink} como plataforma de \textit{Model-Based Design}.
A obra propõe um fluxo de trabalho no qual a lógica de controle é modelada graficamente, validada por simulação digital contra um modelo da planta --- atuando como um gêmeo digital --- e, em seguida, convertida automaticamente para blocos funcionais compatíveis com a IEC 61131-3 via \textit{Simulink PLC Coder} \cite{mathworks:2025}.
Os autores comprovam que a documentação de validação e o código \textit{PLC} podem ser gerados automaticamente a partir do modelo do sistema, reduzindo significativamente o tempo de comissionamento.
A metodologia de \textit{Model-Based Design} apresentada por Duggan et al. \cite{duggan:2022} é diretamente aplicável ao presente projeto como estratégia de verificação: a lógica de controle modelada em \textit{Simulink} serve como referência para comparação com a execução na \textit{VM} embarcada.

\subsection{Juliet \& Romeo: Linguagem e VM para Automação Moderna}
\label{sub:cognibotics}

De forma análoga à solução aqui proposta, a plataforma \textit{Juliet \& Romeo}, desenvolvida pela \textit{Cognibotics} \cite{cognibotics:2025}, representa uma abordagem emergente que combina uma linguagem de programação expressiva (\textit{Juliet}) com uma máquina virtual de tempo real (\textit{Romeo}).
O sistema foi projetado para automação moderna e programação de movimentos, oferecendo garantias de tempo real e comportamento determinístico.
Diferentemente da abordagem do presente projeto, que se concentra na interpretação de \textit{bytecodes Ladder} em microcontroladores de recursos limitados, o \textit{Juliet \& Romeo} opera em plataformas de maior capacidade computacional e visa à coexistência com ambientes IEC 61131-3 existentes, como o \textit{CODESYS} \cite{codesys:2025}, buscando atuar como extensão para funcionalidades avançadas de robótica e inteligência artificial.
Não obstante, o conceito de separação entre linguagem de alto nível e máquina virtual de execução determinística é análogo ao proposto neste trabalho, ainda que a plataforma da \textit{Cognibotics} divirja em abordagem, foco e cenário de aplicação.


%% ------------------------------------------------------------------------- %%
\section{Estado da Prática}
\label{sec:estado-pratica}

O estado da prática em soluções reais para execução de lógicas de controle IEC 61131-3 compreende tanto produtos comerciais consolidados quanto projetos de código aberto.
Esta seção descreve as principais soluções disponíveis no mercado e na comunidade, apontando suas características, pontos positivos e limitações.

\subsection{OpenPLC}
\label{sub:estado-openplc}

O \textit{OpenPLC} \cite{openplc:2025} é a solução de código aberto mais madura para implementação de CLPs em \textit{hardware} de propósito geral.
Compatível com as cinco linguagens da IEC 61131-3, suporta execução em plataformas como \textit{Raspberry Pi}, \textit{Arduino} e PCs industriais com \textit{Linux}.
Sua arquitetura baseia-se no compilador \textit{MatIEC}, que traduz o código IEC 61131-3 para C++ executado nativamente na plataforma embarcada selecionada, contando ainda com servidor \textit{web} para \textit{upload} e gerenciamento de programas, suporte ao protocolo \textit{Modbus} e uma camada de \textit{hardware} configurável.
As limitações do \textit{OpenPLC} residem na compilação estática, que gera um binário único no qual a lógica do usuário é fundida ao \textit{firmware}, exigindo a regravação completa da lógica \textit{Ladder} a cada alteração e impossibilitando atualizações dinâmicas; além disso, a portabilidade para novos \textit{hardwares} exige modificações na camada de \textit{drivers} em baixo nível.

\subsection{CODESYS Runtime}
\label{sub:estado-codesys}

O \textit{CODESYS Runtime} \cite{codesys:2025}, desenvolvido pela empresa alemã \textit{3S-Smart Software Solutions GmbH}, é o principal sistema de desenvolvimento IEC 61131-3 independente de fabricante, presente em mais de mil tipos diferentes de dispositivos de mais de quinhentos fabricantes.
O sistema oferece variantes para PCs industriais com \textit{Windows}, dispositivos \textit{Linux} com processadores \textit{ARM} e controladores virtuais (\textit{vPLCs}) executados em contêineres, gerando código nativo --- e não interpretado --- para uma ampla gama de processadores, o que garante desempenho superior.
Por outro lado, o \textit{CODESYS} é um produto comercial de licenciamento pago, representando uma barreira de custo significativa para projetos acadêmicos e aplicações de pequeno porte; seu código-fonte é proprietário e fechado, impedindo a inspeção e a modificação por pesquisadores, e a adaptação para novos \textit{hardwares} requer a aquisição de um \textit{Runtime Toolkit} proprietário.

\subsection{Simulink PLC Coder}
\label{sub:estado-simulink}

O \textit{Simulink PLC Coder} \cite{mathworks:2025}, extensão da \textit{MathWorks} para o \textit{MATLAB}, gera automaticamente código IEC 61131-3 a partir de modelos \textit{Simulink}, gráficos \textit{Stateflow} e funções \textit{MATLAB}.
O código gerado é exportado em formatos compatíveis com os principais IDEs industriais, incluindo \textit{CODESYS}, \textit{Studio 5000}, \textit{TIA Portal} e \textit{Sysmac Studio}, produzindo também bancos de teste para verificação da correspondência entre a simulação e a execução no CLP físico, além de relatórios de geração de código com rastreabilidade bidirecional entre modelo e código.
A principal limitação dessa solução é pressupor a existência de um CLP alvo convencional ou \textit{Soft PLC} para a execução do código gerado, não oferecendo uma camada de virtualização própria para microcontroladores de propósito geral; adicionalmente, as licenças \textit{MATLAB}/\textit{Simulink} e do \textit{PLC Coder} representam um investimento financeiro significativo.

\subsection{MicroPython}
\label{sub:estado-micropython}

Uma abordagem alternativa observada na prática consiste na utilização de interpretadores de linguagens de \textit{script} genéricas, como o \textit{MicroPython} \cite{micropython:2025}, em microcontroladores para implementação de lógicas de controle simples.
O \textit{MicroPython} executa \textit{bytecodes Python} em uma máquina virtual compatível com plataformas como \textit{ESP32}, \textit{STM32} e \textit{Raspberry Pi Pico}.
Embora ofereça portabilidade e facilidade de programação, a gestão de memória por \textit{Garbage Collection} introduz pausas não determinísticas que violam os requisitos de temporização estrita do ciclo de varredura; além disso, a linguagem \textit{Python} não foi projetada para expressar nativamente primitivas de lógica \textit{Ladder}.
Dessa forma, o \textit{MicroPython} atende satisfatoriamente aplicações de \textit{IoT} e prototipagem, mas não oferece as garantias de determinismo temporal necessárias para controle industrial.


%% ------------------------------------------------------------------------- %%
\section{Conceitos Fundamentais de Arquitetura}
\label{sec:conceitos-arquitetura}

Esta seção introduz os conceitos teóricos que fundamentam as decisões de projeto detalhadas nos capítulos de desenvolvimento deste trabalho.

Uma Arquitetura de Conjunto de Instruções (\textit{ISA}) de largura fixa define instruções que ocupam sempre o mesmo número de \textit{bits} em memória.
Essa escolha simplifica o estágio de busca e decodificação (\textit{fetch-decode}) de um interpretador: como o avanço do ponteiro de instrução é constante, o tempo de decodificação torna-se previsível, característica desejável para sistemas com requisitos de determinismo temporal, como o ciclo de varredura de um CLP.
O ciclo \textit{fetch-decode-execute} constitui o modelo clássico de operação de uma máquina virtual ou processador: a instrução é buscada na memória de programa, decodificada para identificar a operação e os operandos, e então executada, repetindo-se o processo para a instrução seguinte.

A Camada de Abstração de \textit{Hardware} (\textit{HAL}) é um padrão arquitetural que encapsula o acesso aos periféricos físicos de um dispositivo --- como \textit{GPIOs}, temporizadores e conversores analógico-digitais --- por trás de uma interface estável e independente do \textit{hardware} subjacente.
A adoção de uma \textit{HAL} permite que a lógica de aplicação e o motor de execução permaneçam inalterados quando o projeto é portado para uma nova plataforma, restringindo o esforço de adaptação à reimplementação da própria camada.

Um ambiente de execução isolado (\textit{sandbox}) é um mecanismo que restringe o escopo de atuação de um programa em execução, validando em tempo real os limites de seus acessos --- como endereços de memória e índices de entrada e saída --- e rejeitando operações inválidas.
Em sistemas embarcados críticos, esse mecanismo contém erros estruturais na lógica do usuário antes que eles comprometam a estabilidade do restante do sistema.

Por fim, o \textit{Model-Based Design} é uma metodologia de desenvolvimento na qual um modelo executável do sistema --- e, quando aplicável, da planta controlada --- é utilizado como artefato central ao longo de todo o ciclo de vida do projeto, desde a especificação até a verificação.
Nesse paradigma, a simulação do modelo serve como referência comportamental para validar a implementação final, estratégia adotada neste trabalho para a validação cruzada entre a lógica de controle modelada em \textit{Simulink}/\textit{Stateflow} e sua execução na máquina virtual embarcada.

% !TeX root = ../main.tex
%% ------------------------------------------------------------------------- %%
\chapter{Implementação}
\label{cap:implementacao}

Os conceitos que sustentam a solução --- Arquitetura de Conjunto de Instruções (\textit{ISA}), Camada de Abstração de \textit{Hardware} (\textit{HAL}), ambiente de execução isolado (\textit{sandbox}) e \textit{Model-Based Design} --- foram apresentados no Capítulo~\ref{cap:fundamentacao}.
Este capítulo descreve como cada um deles se materializa na implementação: a arquitetura em camadas (Seção~\ref{sec:visao-geral-implementacao}), a plataforma-alvo (Seção~\ref{sec:plataforma-alvo}), a máquina virtual e o conjunto de instruções (Seção~\ref{sec:vm-isa}), o modelo de memória e o ambiente de execução isolado (Seção~\ref{sec:memoria-sandbox}), a camada de abstração de \textit{hardware} (Seção~\ref{sec:hal-implementacao}), o protocolo de comunicação (Seção~\ref{sec:protocolo}), a modelagem em \textit{Simulink}/\textit{Stateflow} (Seção~\ref{sec:stateflow-mbd}) e a estratégia de validação (Seção~\ref{sec:validacao-testes}).
A especificação completa do conjunto de instruções e os ensaios do bloco de controle PID constam, respectivamente, do \ref{ape:isa} e do \ref{ape:pid}.
Cada escolha de projeto é acompanhada de sua justificativa, reunidas no Quadro~\ref{qua:escolhas-projeto}.

\tikzset{
    bloco/.style={draw, rounded corners=2pt, align=center, font=\small, fill=gray!10, inner sep=4pt, minimum height=8mm},
    estado/.style={draw, rounded corners=6pt, align=center, font=\small, fill=gray!10, inner sep=4pt, minimum height=9mm, minimum width=2.2cm},
    seta/.style={-{Stealth[length=2.4mm]}, thick},
    rotulo/.style={font=\footnotesize, inner sep=1.5pt, fill=white}
}

%% ------------------------------------------------------------------------- %%
\section{Visão Geral da Arquitetura Implementada}
\label{sec:visao-geral-implementacao}

A solução é organizada em quatro módulos sobrepostos a uma fronteira de portabilidade, como mostra a Figura~\ref{fig:arquitetura-geral}.
O módulo de máquina virtual reúne o motor de execução e o conjunto de instruções; o módulo de mapas de memória guarda o estado de todos os elementos endereçáveis pelo programa (entradas, saídas, memórias auxiliares, temporizadores, contadores, blocos de controle e registradores de dados); o módulo de comunicação implementa o protocolo com a ferramenta de \textit{host}; e a camada de abstração de \textit{hardware} concentra todo o acesso aos periféricos do microcontrolador.
O programa do usuário é compilado para \textit{bytecode}, transferido ao dispositivo pelo protocolo serial e armazenado como dado somente leitura, de modo que a lógica de aplicação nunca acessa periféricos diretamente: ela é apenas percorrida pela máquina virtual, e esta, por sua vez, só interage com o \textit{hardware} através do contrato da camada de abstração.
Essa fronteira única entre o motor de execução e o microcontrolador é o que torna o motor de execução independente da plataforma, conforme discutido na Seção~\ref{sec:fundamentos-solucao}.

\begin{figure}[htb]
    \centering
    \begin{tikzpicture}[x=1cm,y=1cm]
        \node[bloco, fill=corApp, minimum width=10.4cm] (app) at (0,5.0) {Lógica de aplicação: \textit{bytecode} do usuário (dado somente leitura)};
        \node[bloco, fill=corVM, minimum width=3.0cm, minimum height=14mm, font=\footnotesize] (vm) at (-3.8,3.0) {Máquina virtual:\\motor de execução\\e instruções};
        \node[bloco, fill=corMem, minimum width=3.0cm, minimum height=14mm, font=\footnotesize] (map) at (0,3.0) {Mapas de memória:\\imagens de E/S\\e blocos funcionais};
        \node[bloco, fill=corCom, minimum width=3.0cm, minimum height=14mm, font=\footnotesize] (com) at (3.8,3.0) {Comunicação:\\enquadramento COBS\\e CRC16};
        \node[bloco, fill=corHAL, minimum width=10.4cm] (hal) at (0,1.1) {Camada de abstração de \textit{hardware} (contrato único com o microcontrolador)};
        \node[bloco, fill=corHW, minimum width=10.4cm] (mcu) at (0,-0.1) {Microcontrolador STM32F767ZI e periféricos};
        \node[bloco, fill=white, minimum width=2.4cm, minimum height=14mm] (host) at (8.6,3.0) {Ferramenta\\de \textit{host}};
        \draw[seta] (vm.north |- app.south) -- (vm.north);
        \draw[seta, <->] (vm) -- (map);
        \draw[seta, <-] (com.north |- app.south) -- node[rotulo, right] {carga} (com.north);
        \draw[seta] (vm.south) -- (vm.south |- hal.north);
        \draw[seta] (map.south) -- (map.south |- hal.north);
        \draw[seta] (com.south) -- (com.south |- hal.north);
        \draw[seta] (hal) -- (mcu);
        \draw[seta, <->] (com) -- node[rotulo, above] {serial} (host);
    \end{tikzpicture}
    \caption{Visão geral dos módulos do firmware e fronteira de portabilidade}
    \label{fig:arquitetura-geral}

    \fonte{Elaborado pelo autor (2026).}
\end{figure}

O laço principal do firmware repete indefinidamente quatro etapas: a leitura das entradas físicas com filtro de \textit{debounce} configurável, a execução de uma varredura completa do \textit{bytecode} pela máquina virtual, a aplicação da imagem de saídas ao \textit{hardware} e o atendimento da comunicação.
Esse ciclo de varredura (\textit{scan cycle}) reproduz, em software, o funcionamento cíclico de um CLP descrito na Seção~\ref{sec:tema} \cite{bolton:2015}.
O Quadro~\ref{qua:recursos} reúne a capacidade dos recursos disponíveis ao programa do usuário, todos definidos por constantes de compilação e, portanto, ajustáveis sem alteração do motor de execução.

\begin{quadro}[htb]
    \centering
    \begin{tabular}{|l|c|p{7.2cm}|}
        \hline
        Recurso & Quantidade & Observação \\
        \hline
        Entradas digitais & 10 & Com filtro de \textit{debounce} \\
        \hline
        Saídas digitais & 12 & --- \\
        \hline
        Entradas analógicas & 8 & Conversor de 12 \textit{bits} \\
        \hline
        Saídas PWM & 4 & Uma delas analógica pura (0 a 10 V) \\
        \hline
        Temporizadores & 32 & TON, TOF e TP; bases de 10 ms a 1 min \\
        \hline
        Contadores & 32 & CTU, CTD e CTUD; valor de 32 \textit{bits} com sinal \\
        \hline
        Blocos PID e RAMP & 16 + 16 & Instâncias independentes \\
        \hline
        Pilha lógica e pilha de blocos & 16 + 16 & Ramificações e blocos aninhados \\
        \hline
    \end{tabular}

    \caption{Recursos disponíveis ao programa do usuário}
    \label{qua:recursos}
    \fonte{Elaborado pelo autor (2026).}
\end{quadro}


%% ------------------------------------------------------------------------- %%
\section{Plataforma-Alvo}
\label{sec:plataforma-alvo}

A implementação tem como alvo o microcontrolador STM32F767ZI, de núcleo \textit{Arm Cortex-M7} com frequência de até 216 MHz e unidade de ponto flutuante \cite{st:ds11532}.
Quatro requisitos da solução justificam essa escolha.
O primeiro é a aritmética de ponto flutuante: 13 das instruções aritméticas e 6 das de comparação operam sobre \textit{F32}, e a unidade de ponto flutuante do núcleo as executa em \textit{hardware}, com tempo previsível e sem bibliotecas de emulação, o que preserva o determinismo discutido na Seção~\ref{sec:problema-literatura} \cite{liu:2000}.
O segundo é a memória: o dispositivo oferece até 2 MB de \textit{Flash} e 512 KB de SRAM, incluindo memórias acopladas ao núcleo (128 KB de dados e 16 KB de instruções) destinadas a dados e rotinas críticas de tempo real \cite{st:ds11532}, o que permite alocar os mapas de memória e o laço de varredura em regiões de acesso previsível.
O terceiro é a cobertura de periféricos: conversores analógico-digitais de 12 \textit{bits}, temporizadores, controladores de DMA, interfaces seriais e unidade de cálculo de CRC atendem a todos os grupos de funções do contrato da camada de abstração (Seção~\ref{sec:hal-implementacao}) \cite{st:ds11532,st:rm0410}.
O quarto é o ferramental: a placa NUCLEO-F767ZI integra o depurador/gravador ST-LINK, dispensando sonda externa e disponibilizando um canal serial para o \textit{host} \cite{st:um1974}, e a ferramenta \textit{STM32CubeMX} gera a inicialização dos periféricos a partir de uma descrição gráfica.
Como o desempenho do núcleo é superior ao necessário para o interpretador, espera-se que o tempo de varredura seja determinado pelo tamanho do programa do usuário e não pelo custo de interpretação, hipótese verificada pelos ensaios de desempenho descritos na Seção~\ref{sec:validacao-testes}.

A Figura~\ref{fig:esboco-placa} apresenta o esboço, em blocos, da conexão entre a placa e o circuito de entradas e saídas.
A atribuição individual de pinos é definida no projeto do circuito e não altera o contrato da camada de abstração.

\begin{figure}[htb]
    \centering
    \begin{tikzpicture}[x=1cm,y=1cm]
        \node[bloco, fill=corHW, minimum width=3.6cm, minimum height=34mm] (mcu) at (0,0) {\textbf{NUCLEO-F767ZI}\\STM32F767ZI\\[2pt]GPIO\quad ADC\\TIM (PWM)\\USART + DMA\\CRC};
        \node[bloco, fill=corCom, minimum width=3.0cm] (di) at (-5.6,1.1) {10 entradas\\digitais};
        \node[bloco, fill=corCom, minimum width=2.3cm] (dic) at (-5.6,-0.3) {Condicionamento\\de sinais};
        \node[bloco, fill=corCom, minimum width=3.0cm] (ai) at (-5.6,-1.7) {8 entradas\\analógicas};
        \node[bloco, fill=corMem, minimum width=3.0cm] (do) at (5.6,1.5) {12 saídas\\digitais};
        \node[bloco, fill=corMem, minimum width=3.0cm] (pwm) at (5.6,0.1) {3 saídas PWM};
        \node[bloco, fill=corMem, minimum width=3.0cm] (ao) at (5.6,-1.4) {1 saída analógica\\0 a 10 V};
        \node[bloco, fill=white, minimum width=3.6cm] (pc) at (0,-3.4) {Ferramenta de \textit{host}\\(PC)};
        \draw[seta] (di.east) -- ++(0.5,0) |- ([yshift=6mm]mcu.west);
        \draw[seta] (ai.east) -- ++(0.5,0) |- ([yshift=-8mm]mcu.west);
        \draw[seta] ([yshift=8mm]mcu.east) -- ++(0.5,0) |- (do.west);
        \draw[seta] (mcu.east) -- ++(0.5,0) |- (pwm.west);
        \draw[seta] ([yshift=-8mm]mcu.east) -- ++(0.5,0) |- (ao.west);
        \draw[seta, <->] (mcu.south) -- node[rotulo, right] {porta serial virtual (ST-LINK)} (pc.north);
    \end{tikzpicture}
    \caption{Esboço em blocos da conexão entre a placa-alvo e o circuito de entradas e saídas}
    \label{fig:esboco-placa}

    \fonte{Elaborado pelo autor (2026), com base em \cite{st:um1974}.}
\end{figure}


%% ------------------------------------------------------------------------- %%
\section{Máquina Virtual e Conjunto de Instruções}
\label{sec:vm-isa}

\subsection{Ciclo de execução}
\label{sub:ciclo-execucao}

O motor de execução implementa o ciclo \textit{fetch-decode-execute} apresentado na Seção~\ref{sec:fundamentos-solucao}.
A cada varredura, o contador de instruções é posicionado no início do \textit{bytecode}, a pilha lógica e o acumulador são zerados e a imagem de saídas retentivas é reaplicada; em seguida, cada instrução é lida da memória de programa (\textit{fetch}), tem o campo de opcode isolado (\textit{decode}) e é despachada para a rotina correspondente (\textit{execute}) até que a instrução de fim de varredura (\textit{EOS}) seja alcançada, como ilustra a Figura~\ref{fig:fetch-decode-execute}.
Um opcode fora do repertório válido interrompe a varredura e leva a máquina ao estado de erro (\textit{trap}), no qual a execução permanece bloqueada, com um pino de sinalização acionado, até que um novo \textit{bytecode} seja transferido.

\begin{figure}[htb]
    \centering
    \begin{tikzpicture}[x=1cm,y=1cm]
        \node[estado, fill=corHAL, minimum width=2.0cm] (ini) at (-6.6,0) {Início da\\varredura};
        \node[estado, fill=corApp, minimum width=2.0cm] (fetch) at (-3.3,0) {\textit{Fetch}};
        \node[estado, fill=corCom, minimum width=2.0cm] (decode) at (0,0) {\textit{Decode}};
        \node[estado, fill=corMem, minimum width=2.0cm] (exec) at (3.3,0) {\textit{Execute}};
        \node[estado, fill=corHAL, minimum width=2.0cm] (fim) at (6.6,0) {Fim da\\varredura};
        \node[estado, fill=corErr, minimum width=2.0cm] (trap) at (0,-2.4) {Erro (\textit{trap})};
        \draw[seta] (ini) -- (fetch);
        \draw[seta] (fetch) -- (decode);
        \draw[seta] (decode) -- node[rotulo, above=1pt, fill=none] {válido} (exec);
        \draw[seta] (exec) -- node[rotulo, above=1pt, fill=none] {\textit{EOS}} (fim);
        \draw[seta] (exec.north) -- ++(0,0.9) -| node[rotulo, pos=0.25, above] {demais instruções} (fetch.north);
        \draw[seta] (decode.south) -- node[rotulo, right] {inválido} (trap.north);
    \end{tikzpicture}
    \caption{Ciclo \textit{fetch-decode-execute} da máquina virtual}
    \label{fig:fetch-decode-execute}

    \fonte{Elaborado pelo autor (2026).}
\end{figure}

O despacho é feito por uma tabela de 128 posições, uma para cada valor possível do opcode de 7 \textit{bits}, cada uma apontando para a rotina da instrução.
Essa escolha, em vez de uma cadeia de comparações sequenciais, torna o tempo de despacho constante --- acesso a um vetor --- e independente de qual instrução é executada, característica desejável para a previsibilidade do ciclo de varredura \cite{liu:2000,patterson:2020}.
Antes da primeira varredura de cada programa carregado, o \textit{bytecode} é percorrido duas vezes: uma validação rejeita opcodes inválidos, e um pré-processamento extrai, uma única vez, as bases de tempo de temporizadores, blocos PID e blocos RAMP, retirando esse trabalho estático do laço de varredura.

\subsection{Formato e repertório de instruções}
\label{sub:formato-repertorio}

Toda instrução começa por uma palavra-base de 32 \textit{bits}, alinhada em palavra, que contém o opcode e o operando principal.
Essa primeira palavra segue sempre um de dois formatos, mostrados na Figura~\ref{fig:formato-instrucao}: o genérico, usado por instruções que endereçam um operando (tipo, índice e sub-tipo), e o dos blocos PID e RAMP, que identificam a instância e a base de tempo.
Por isso o estágio de decodificação é uniforme, pois lê sempre a mesma quantidade de \textit{bytes} para identificar a instrução \cite{patterson:2020}.
O \textit{bit} mais significativo do \textit{byte} de opcode é reservado à depuração e mascarado antes do despacho.

O tamanho total da instrução, entretanto, não é fixo: a palavra-base pode vir acompanhada de uma estrutura auxiliar, de dois tipos.
O primeiro tipo são palavras adicionais com constantes imediatas de 32 \textit{bits}, colocadas logo após a palavra-base.
Uma operação aritmética, de comparação ou de conversão com operando constante ocupa duas palavras, e o mesmo vale para um temporizador ou contador com valor programado constante.
O escalonamento (\texttt{SCALE}) é o caso extremo: carrega sempre os quatro limites de faixa e, se o valor de entrada também for constante, um quinto valor, ocupando cinco ou seis palavras (20 ou 24 \textit{bytes}).
O número de palavras adicionais é determinado pelo opcode e pelo campo de tipo da própria palavra-base, de modo que é conhecido já na decodificação e é sempre múltiplo de palavra, o que preserva o alinhamento e permite à validação percorrer o programa sem marcadores de fim de instrução.

O segundo tipo de estrutura auxiliar são instruções que preparam o estado de outra.
As instruções \texttt{PID}, \texttt{RAMP} e \texttt{PWM\_OUT} ocupam apenas a palavra-base, e os seus parâmetros (\textit{setpoint}, ganhos, limites, taxas, \textit{duty} e frequência) são gravados antes por instruções \texttt{STW}.
Da mesma forma, as entradas de temporizadores e contadores (habilitação, contagem, \textit{reset} e carga) são gravadas por instruções \texttt{ST} anteriores a \texttt{TIM} e \texttt{COUNT}.
Essa separação mantém a palavra-base pequena e o formato uniforme, ao custo de o comportamento de um bloco funcional depender de instruções que o precedem, o que orienta a composição dos programas de teste dos blocos (Seção~\ref{sec:validacao-testes}).
O \ref{ape:isa} traz o tamanho de cada instrução.

\begin{figure}[htb]
    \centering
    \begin{tikzpicture}[x=1cm,y=1cm, font=\footnotesize]
        \def\bw{0.44}
        \foreach \b/\x in {0/0, 1/3.52, 2/7.04, 3/10.56} {
            \node[font=\footnotesize] at ({\x+1.76},1.95) {byte \b};
        }
        \node[anchor=east] at (-0.2,1.0) {Genérico};
        \tikzset{campo/.style={draw, minimum height=8mm, font=\scriptsize, inner sep=0pt, align=center, anchor=south west}}
        \node[campo, fill=corErr, minimum width=0.44cm] at (0,0.6) {D};
        \node[campo, fill=corVM, minimum width=3.08cm] at (0.44,0.6) {opcode};
        \node[campo, fill=corApp, minimum width=1.76cm] at (3.52,0.6) {tipo};
        \node[campo, fill=corMem, minimum width=1.76cm] at (5.28,0.6) {índice alto};
        \node[campo, fill=corMem, minimum width=1.76cm] at (7.04,0.6) {índice baixo};
        \node[campo, fill=corCom, minimum width=1.76cm] at (8.80,0.6) {sub-tipo};
        \node[campo, fill=corHAL, minimum width=3.52cm] at (10.56,0.6) {reservado};
        \node[anchor=east] at (-0.2,-0.4) {PID e RAMP};
        \node[campo, fill=corErr, minimum width=0.44cm] at (0,-0.8) {D};
        \node[campo, fill=corVM, minimum width=3.08cm] at (0.44,-0.8) {opcode};
        \node[campo, fill=corMem, minimum width=1.76cm] at (3.52,-0.8) {instância};
        \node[campo, fill=corCom, minimum width=0.88cm] at (5.28,-0.8) {TB};
        \node[campo, fill=corHAL, minimum width=0.88cm] at (6.16,-0.8) {res.};
        \node[campo, fill=corHAL, minimum width=3.52cm] at (7.04,-0.8) {reservado};
        \node[campo, fill=corHAL, minimum width=3.52cm] at (10.56,-0.8) {reservado};
        \node[anchor=west, font=\scriptsize, align=left] at (0,-1.5) {D: \textit{bit} de depuração (1); opcode (7); tipo, índices, sub-tipo e instância (4 \textit{bits} cada);\\ TB: base de tempo (2); res.: reservado (2); reservado (8).};
    \end{tikzpicture}
    \caption{Formato da palavra-base de instrução (32 \textit{bits})}
    \label{fig:formato-instrucao}

    \fonte{Elaborado pelo autor (2026).}
\end{figure}

O repertório contém 76 instruções, com opcodes de \texttt{0x00} a \texttt{0x4B}, além da instrução de fim de varredura (\texttt{0x7F}), organizadas nos grupos do Quadro~\ref{qua:grupos-instrucoes}.
A lógica de contatos e bobinas segue o modelo de \textit{Instruction List} da norma IEC 61131-3 \cite{iec:2025}, com pilha lógica para ramificações e pilha de blocos para associações série e paralelo de trechos de \textit{rung}.
A especificação de cada opcode, dos tipos de operando e dos sub-tipos de cada bloco funcional consta do \ref{ape:isa}.

\begin{quadro}[htb]
    \centering
    \begin{tabular}{|p{6.1cm}|c|c|}
        \hline
        Grupo & Instruções & Opcodes \\
        \hline
        Controle de fluxo e pilhas & 7 & \texttt{0x00}, \texttt{0x01}, \texttt{0x12}--\texttt{0x16} \\
        \hline
        Contatos e lógica booleana (incluindo bordas) & 15 & \texttt{0x03}--\texttt{0x11} \\
        \hline
        Armazenamento e bobinas (\texttt{ST}, \texttt{OUT}, \texttt{SET}, \texttt{RESET}, \texttt{TOGGLE}) & 5 & \texttt{0x02}, \texttt{0x17}--\texttt{0x1A} \\
        \hline
        Temporizadores e contadores & 2 & \texttt{0x1B}--\texttt{0x1C} \\
        \hline
        Aritmética e movimentação (\textit{S32}/\textit{F32}) & 26 & \texttt{0x1D}--\texttt{0x36} \\
        \hline
        Comparação (\textit{S32}/\textit{F32}) & 12 & \texttt{0x37}--\texttt{0x42} \\
        \hline
        Conversão e escalonamento & 4 & \texttt{0x43}--\texttt{0x46} \\
        \hline
        Saídas analógica e PWM, escrita em bloco e blocos PID e RAMP & 5 & \texttt{0x47}--\texttt{0x4B} \\
        \hline
        Total & 76 & \texttt{0x00}--\texttt{0x4B} \\
        \hline
    \end{tabular}

    \caption{Grupos de instruções da máquina virtual}
    \label{qua:grupos-instrucoes}
    \fonte{Elaborado pelo autor (2026).}
\end{quadro}

\subsection{Blocos funcionais com estado}
\label{sub:blocos-funcionais}

Blocos funcionais mantêm estado entre varreduras e, por isso, não o armazenam no \textit{bytecode}, que é somente leitura, mas em mapas de memória dedicados (Seção~\ref{sec:memoria-sandbox}).
Os temporizadores implementam os três tipos da IEC 61131-3 --- \textit{TON} (retardo na energização), \textit{TOF} (retardo na desenergização) e \textit{TP} (pulso) --- com tempo programado constante ou lido de registrador e quatro bases de tempo (10 ms, 100 ms, 1 s e 1 min).
A contagem parte da borda de subida de \texttt{EN} no \textit{TON} e no \textit{TP} (este não retentável durante o pulso) e da borda de descida no \textit{TOF}, e é medida contra a base de tempo do sistema, e não pela contagem de varreduras, para que o intervalo não dependa do tamanho do programa.
Alterações do tempo programado durante a contagem são absorvidas sem reiniciá-la.
Os contadores implementam \textit{CTU}, \textit{CTD} e \textit{CTUD}, com entradas de contagem sensíveis a borda, entradas de \textit{reset} e de carga (\textit{load}) e valor programado constante ou de registrador.
O bloco \textit{RAMP} é um limitador de taxa de variação: sua saída persegue o valor alvo com taxas máximas independentes de subida e de descida, por unidade de tempo configurável; com \texttt{EN} desligado a saída é zero, e na borda de subida de \texttt{EN} ela assume imediatamente o valor inicial, sem rampa.

O bloco PID implementa o algoritmo posicional na forma paralela, com quatro decisões de projeto justificadas pelo comportamento em malha fechada \cite{astrom:2006}.
Primeiro, o intervalo de amostragem é fixo e configurável, e não medido: o bloco só recalcula quando esse período decorre, o que impede que a variação do tempo de varredura entre nos termos integral e derivativo.
Segundo, o termo derivativo é calculado sobre a variável de processo e não sobre o erro, evitando o pico de saída (\textit{derivative kick}) quando o \textit{setpoint} muda.
Terceiro, o anti-\textit{windup} é feito por integração condicional (\textit{clamping}): a integração é suspensa quando a saída já está saturada e o erro a empurra mais fundo na mesma direção, e retomada assim que o erro inverte de sinal.
Quarto, no modo manual o integrador acompanha a saída (\textit{back-tracking}), de modo que a passagem de manual para automático ocorre sem salto (\textit{bumpless transfer}).
A banda morta é aplicada ao erro antes dos termos proporcional e integral, e a saída é sempre limitada ao intervalo configurado, inclusive no modo manual.


%% ------------------------------------------------------------------------- %%
\section{Modelo de Memória e Ambiente de Execução Isolado}
\label{sec:memoria-sandbox}

Toda a memória do sistema é alocada estaticamente em tempo de compilação: não há alocação dinâmica nem coleta de lixo, o que elimina fragmentação e a variação de tempo que esses mecanismos introduziriam \cite{liu:2000}.
Cada tipo de elemento endereçável possui um mapa dedicado --- imagens de entradas, saídas e memórias auxiliares (\textit{bits}), registradores de dados de 32 \textit{bits}, temporizadores, contadores, instâncias PID e RAMP, e entradas e saídas analógicas ---, e o programa do usuário os referencia apenas por tipo e índice, nunca por ponteiro, como mostra a Figura~\ref{fig:mapa-memoria}.

\begin{figure}[htb]
    \centering
    \begin{tikzpicture}[x=1cm,y=1cm]
        \node[bloco, fill=corVM, minimum width=3.6cm, minimum height=60mm, font=\small] (vm) at (0,0) {\textbf{Máquina virtual}\\[6pt]\footnotesize acesso por tipo\\e índice, validado\\a cada operação};
        \node[bloco, fill=corApp, minimum width=6.8cm, minimum height=9mm] (bc) at (7.2,2.4) {\textit{Bytecode} (somente leitura)};
        \node[bloco, fill=corMem, minimum width=6.8cm, minimum height=9mm] (img) at (7.2,1.2) {Imagens de entradas, saídas e memórias (\textit{bits})};
        \node[bloco, fill=corHW, minimum width=6.8cm, minimum height=9mm] (dr) at (7.2,0) {Registradores de dados (32 \textit{bits})};
        \node[bloco, fill=corRos, minimum width=6.8cm, minimum height=9mm] (fb) at (7.2,-1.2) {Temporizadores, contadores, PID e RAMP};
        \node[bloco, fill=corCom, minimum width=6.8cm, minimum height=9mm] (an) at (7.2,-2.4) {Entradas analógicas e saídas PWM};
        \draw[seta, <-] (vm.east |- bc.west) -- (bc.west);
        \draw[seta, <->] (vm.east |- img.west) -- (img.west);
        \draw[seta, <->] (vm.east |- dr.west) -- (dr.west);
        \draw[seta, <->] (vm.east |- fb.west) -- (fb.west);
        \draw[seta, <->] (vm.east |- an.west) -- (an.west);
    \end{tikzpicture}
    \caption{Regiões de memória e forma de acesso pela máquina virtual}
    \label{fig:mapa-memoria}

    \fonte{Elaborado pelo autor (2026).}
\end{figure}

Essa segregação materializa o ambiente de execução isolado descrito na Seção~\ref{sec:fundamentos-solucao}.
O \textit{bytecode} é rejeitado se contiver opcode inválido, e as funções de acesso de cada mapa validam o índice a cada operação, e uma leitura fora da faixa retorna um valor neutro (falso ou zero) sem acessar memória, de modo que um erro de endereçamento na lógica do usuário não alcança memória fora das regiões alocadas nem o estado interno da máquina.
Diferentemente da compilação para código nativo, na qual um erro de lógica executa diretamente no processador, aqui o pior efeito possível é local à instrução defeituosa ou, no caso de opcode inválido, a parada controlada da execução com sinalização externa.


%% ------------------------------------------------------------------------- %%
\section{Camada de Abstração de Hardware}
\label{sec:hal-implementacao}

O contrato da camada de abstração expõe somente funções e tipos primitivos, organizados por grupo de periférico, e nenhum tipo do fabricante do microcontrolador aparece nessa interface: eles ficam confinados à implementação.
É essa separação entre contrato e implementação que faz da camada o único ponto do firmware dependente do microcontrolador, e que reduz a adaptação a outra plataforma à reimplementação dela.
O Quadro~\ref{qua:hal} associa cada grupo de funções ao periférico do STM32F767ZI que o implementa.

\begin{quadro}[htb]
    \centering
    \begin{tabular}{|p{5.6cm}|p{8.6cm}|}
        \hline
        Grupo de funções & Periférico e observações \\
        \hline
        Entradas e saídas digitais; sinalização de erro & \textit{GPIO} \\
        \hline
        Comunicação serial com o \textit{host} & \textit{USART} com DMA em transmissão e recepção, sobre a porta serial virtual do ST-LINK \\
        \hline
        Base de tempo (ms e $\mu$s) e temporizador de varredura & \textit{Timer} de propósito geral \\
        \hline
        Entradas analógicas (8 canais) & ADC de 12 \textit{bits} \\
        \hline
        Saídas PWM (4 canais) e saída analógica & \textit{Timer} em modo PWM; filtro externo para a saída de 0 a 10 V \\
        \hline
        Verificação de integridade (CRC16) & Unidade de CRC do microcontrolador ou rotina em \textit{software}, selecionada na inicialização \\
        \hline
    \end{tabular}

    \caption{Grupos de funções da camada de abstração e periféricos do STM32F767ZI}
    \label{qua:hal}
    \fonte{Elaborado pelo autor (2026), com base em \cite{st:ds11532,st:rm0410,st:um1974}.}
\end{quadro}

A implementação utiliza a biblioteca de \textit{HAL} fornecida pelo fabricante (\textit{STM32Cube HAL}), em vez do acesso direto a registradores (\textit{Low-Layer}), por dois motivos.
O código de inicialização gerado pela ferramenta \textit{STM32CubeMX} reduz o risco de erro de configuração de registradores, cuja depuração consome tempo desproporcional ao ganho de desempenho da alternativa de baixo nível nesta aplicação; e, como a camada de abstração do projeto já é uma fronteira própria, o eventual sobrecusto da biblioteca do fabricante fica contido nela e não se propaga ao motor de execução.
Interfaces de armazenamento externo e de USB não fazem parte do escopo desta monografia.


%% ------------------------------------------------------------------------- %%
\section{Protocolo de Comunicação}
\label{sec:protocolo}

O protocolo transfere o \textit{bytecode} do \textit{host} ao dispositivo e permite a depuração em tempo real (observação e forçamento de variáveis, pausa e retomada da execução).
Cada mensagem é um quadro composto pelo comando, um identificador de \textit{hardware}, os dados e uma verificação de integridade CRC16 de polinômio 0x8005 (forma refletida 0xA001), e é enviada com enquadramento \textit{Consistent Overhead Byte Stuffing} (\textit{COBS}).
O \textit{COBS} elimina o byte zero do conteúdo do quadro com sobrecarga limitada e conhecida, o que permite usar esse byte como delimitador sem ambiguidade \cite{cheshire:1999}; o CRC16 detecta corrupção em canais seriais ruidosos.
O identificador de dispositivo presente em cada quadro permite ao \textit{firmware} rejeitar comandos endereçados a um dispositivo diferente do esperado.
Os pacotes de dados têm no máximo 512 \textit{bytes}, dos quais até 256 de carga útil.

Os comandos organizam-se em três grupos: controle (confirmação positiva, confirmação negativa com código de erro e notificação de erro da máquina virtual), transferência (início, pacote de dados e fim) e depuração (início e fim, pausa e retomada, publicação e remoção de variáveis observadas, forçamento de valores e escrita direta em memória e registradores).
O dispositivo mantém uma máquina de estados com três modos --- espera, transferência e depuração ---, e em cada um deles rejeita, com um código de \textit{NACK} específico, qualquer comando inválido, de tamanho inconsistente, com CRC incorreto ou destinado a outro identificador de \textit{hardware}.
A Figura~\ref{fig:protocolo-sequencia} apresenta a sequência de uma transferência de \textit{bytecode}.
A autenticação do \textit{host} não faz parte do escopo desta monografia; a rejeição correta de quadros malformados é, portanto, a barreira de robustez do canal contra ruído e contra uso incorreto da ferramenta.

\begin{figure}[htb]
    \centering
    \begin{tikzpicture}[x=1cm,y=1cm]
        \node[bloco, fill=corApp, minimum width=3cm] (h) at (-4,0) {Ferramenta de \textit{host}};
        \node[bloco, fill=corVM, minimum width=3cm] (d) at (4,0) {Dispositivo};
        \draw[thick] (h.south) -- (-4,-6.4);
        \draw[thick] (d.south) -- (4,-6.4);
        \draw[seta] (-4,-1.0) -- node[rotulo, above] {\texttt{START\_TRANSFER}} (4,-1.0);
        \draw[seta, dashed] (4,-1.7) -- node[rotulo, above] {\texttt{ACK} (CRC e identificador válidos) ou \texttt{NACK}} (-4,-1.7);
        \draw[seta] (-4,-2.7) -- node[rotulo, above] {\texttt{DATA\_PACKET} $n$ (COBS + CRC16)} (4,-2.7);
        \draw[seta, dashed] (4,-3.4) -- node[rotulo, above] {\texttt{ACK} do pacote $n$ ou \texttt{NACK}} (-4,-3.4);
        \node[rotulo, fill=none] at (0,-4.05) {$\vdots$ repetido para cada pacote};
        \draw[seta] (-4,-4.8) -- node[rotulo, above] {\texttt{STOP\_TRANSFER}} (4,-4.8);
        \draw[seta, dashed] (4,-5.5) -- node[rotulo, above] {\texttt{ACK} ou \texttt{NACK}} (-4,-5.5);
    \end{tikzpicture}
    \caption{Sequência de uma transferência de \textit{bytecode}}
    \label{fig:protocolo-sequencia}

    \fonte{Elaborado pelo autor (2026).}
\end{figure}


%% ------------------------------------------------------------------------- %%
\section{Modelagem em Simulink/Stateflow segundo Model-Based Design}
\label{sec:stateflow-mbd}

A metodologia de \textit{Model-Based Design}, fundamentada na Seção~\ref{sec:fundamentos-solucao}, prescreve um modelo executável do comportamento esperado como referência para validar a implementação \cite{nicolescu:2009}.
Neste trabalho, os modelos de referência cumprem três papéis: servem de oráculo para os ensaios de lógica de contatos, temporização e contagem (Seção~\ref{sec:validacao-testes}), à maneira do modelo semântico executável empregado por Sadolewski e Trybus \cite{sadolewski:2024}; especificam a máquina de modos do bloco PID; e, em conjunto com modelos de planta, fecham a malha dos ensaios \textit{Hardware-in-the-Loop} do \ref{ape:pid}.
O procedimento de validação cruzada, análogo ao de Duggan et al. \cite{duggan:2022}, aplica o mesmo vetor de estímulos ao modelo simulado e à máquina virtual em execução, registra a imagem de saídas de cada varredura e compara os dois resultados bit a bit, com as transições de estado --- e não apenas o valor final --- incluídas na comparação.

A máquina de modos do bloco PID foi escolhida como primeiro artefato de modelagem por estar inteiramente especificada por estados e transições: quatro modos com prioridade estrita, avaliados nesta ordem a cada execução do bloco --- \textsc{Desabilitado}, \textsc{Reset}, \textsc{Manual} e \textsc{Automático} ---, como mostram a Figura~\ref{fig:stateflow-pid} e o Quadro~\ref{qua:estados-pid}.

\begin{figure}[htb]
    \centering
    \begin{tikzpicture}[x=1cm,y=1cm]
        \node[diamond, draw, aspect=2.2, align=center, font=\footnotesize, inner sep=1pt, fill=corHW] (en) at (0,0) {\texttt{EN}?};
        \node[diamond, draw, aspect=2.2, align=center, font=\footnotesize, inner sep=1pt, fill=corHW] (rs) at (4.6,0) {\texttt{RESET}?};
        \node[diamond, draw, aspect=2.2, align=center, font=\footnotesize, inner sep=1pt, fill=corHW] (mn) at (9.2,0) {\texttt{MANUAL}?};
        \node[estado, fill=corHAL] (des) at (0,-2.6) {\textsc{Desabilitado}};
        \node[estado, fill=corErr] (res) at (4.6,-2.6) {\textsc{Reset}};
        \node[estado, fill=corCom] (man) at (9.2,-2.6) {\textsc{Manual}};
        \node[estado, fill=corMem] (aut) at (12.6,0) {\textsc{Automático}};
        \draw[seta] (en) -- node[rotulo, above] {sim} (rs);
        \draw[seta] (rs) -- node[rotulo, above] {não} (mn);
        \draw[seta] (mn) -- node[rotulo, above] {não} (aut);
        \draw[seta] (en) -- node[rotulo, left] {não} (des);
        \draw[seta] (rs) -- node[rotulo, left] {sim} (res);
        \draw[seta] (mn) -- node[rotulo, left] {sim} (man);
    \end{tikzpicture}
    \caption{Prioridade dos modos do bloco PID (especificação do \textit{chart} \textit{Stateflow})}
    \label{fig:stateflow-pid}

    \fonte{Elaborado pelo autor (2026).}
\end{figure}

\begin{quadro}[htb]
    \centering
    \begin{tabular}{|l|p{4.6cm}|p{7.2cm}|}
        \hline
        Modo & Condição & Comportamento \\
        \hline
        \textsc{Desabilitado} & \texttt{EN} = falso & Bloco não calcula; \texttt{CV} assume o limite mínimo de saída; o integrador é preservado, para retomada próxima do ponto de parada. \\
        \hline
        \textsc{Reset} & \texttt{EN} verdadeiro e \texttt{RESET} verdadeiro & Integrador zerado; \texttt{CV} igual ao limite mínimo. \\
        \hline
        \textsc{Manual} & \texttt{EN} verdadeiro, \texttt{RESET} falso e \texttt{MANUAL} verdadeiro & \texttt{CV} segue o valor manual, limitado à faixa; integrador em \textit{back-tracking} para transição sem salto. \\
        \hline
        \textsc{Automático} & \texttt{EN} verdadeiro, \texttt{RESET} e \texttt{MANUAL} falsos & PID posicional com banda morta e anti-\textit{windup} por integração condicional. \\
        \hline
    \end{tabular}

    \caption{Especificação da máquina de modos do bloco PID}
    \label{qua:estados-pid}
    \fonte{Elaborado pelo autor (2026).}
\end{quadro}


%% ------------------------------------------------------------------------- %%
\section{Estratégia de Validação e Testes}
\label{sec:validacao-testes}

A validação é organizada em níveis de teste, escolhidos de acordo com o que cada um comprova isoladamente, conforme a prática da engenharia de \textit{software} descrita na norma ISO/IEC/IEEE 29119 \cite{iso:29119-1}.
Os \textbf{testes de unidade} isolam um grupo de instruções e comparam a imagem de saídas produzida com a de um oráculo.
Os \textbf{testes de integração} exercitam o protocolo de comunicação, o carregamento do \textit{bytecode} e a execução conjuntamente, inclusive com entradas malformadas.
Os \textbf{testes de sistema} fecham a malha de controle com uma planta simulada em tempo real (\textit{Hardware-in-the-Loop}, \textit{HIL}), técnica que permite validar o controlador real sem expor um processo físico a falhas \cite{isermann:1999}; sua aplicação ao bloco PID é apresentada no \ref{ape:pid}.
Os \textbf{testes de desempenho} medem o tempo de varredura e o consumo de memória, atendendo ao objetivo de análise quantitativa da solução.

A infraestrutura de teste é composta por quatro elementos.
Programas de teste são gerados por uma biblioteca de montagem de \textit{bytecode} em \textit{Python}, com uma função por instrução que codifica a palavra-base e as palavras auxiliares, o que dispensa o cálculo manual dos campos.
\textit{Scripts} no \textit{host} transferem o \textit{bytecode}, forçam valores de entrada e observam a imagem de saídas por meio dos comandos de depuração do protocolo.
Modelos em \textit{Simulink}/\textit{Stateflow} e cálculos em \textit{MATLAB} fornecem o resultado esperado (oráculo), e diagramas de temporização da própria norma IEC 61131-3 servem de referência para temporizadores, contadores e bordas \cite{iec:2025}.
Por fim, um instrumento de aquisição analógica (\textit{Analog Discovery 3}) mede as saídas analógicas e fecha a malha nos ensaios \textit{HIL}.
A Figura~\ref{fig:fluxo-validacao} resume o fluxo de um caso de teste.

\begin{figure}[htb]
    \centering
    \begin{tikzpicture}[x=1cm,y=1cm]
        \node[bloco, fill=corApp, minimum width=2.7cm] (prog) at (0,0) {Programa de\\teste (montador)};
        \node[bloco, fill=corCom, minimum width=2.7cm] (down) at (3.6,0) {Transferência\\pelo protocolo};
        \node[bloco, fill=corVM, minimum width=2.7cm] (run) at (7.2,0) {Execução na\\máquina virtual};
        \node[bloco, fill=corMem, minimum width=2.7cm] (obs) at (10.8,0) {Imagem de saídas\\(depuração)};
        \node[bloco, fill=corApp, minimum width=2.7cm] (est) at (0,-2.3) {Vetor de\\estímulos};
        \node[bloco, fill=corRos, minimum width=2.7cm] (orac) at (7.2,-2.3) {Oráculo (modelo,\\cálculo, norma)};
        \node[estado, minimum width=2.7cm, fill=corHW] (cmp) at (10.8,-2.3) {Comparação\\e veredito};
        \draw[seta] (prog) -- (down);
        \draw[seta] (down) -- (run);
        \draw[seta] (run) -- (obs);
        \draw[seta] (obs) -- (cmp);
        \draw[seta] (orac) -- (cmp);
        \draw[seta] (est) -- (orac);
        \draw[seta] (est.north) -- (prog.south);
        \draw[seta] (est.east) -- ++(0.6,0) |- ([yshift=-4mm]run.south west) -- (run.south west);
    \end{tikzpicture}
    \caption{Fluxo de execução e avaliação de um caso de teste}
    \label{fig:fluxo-validacao}

    \fonte{Elaborado pelo autor (2026).}
\end{figure}

O critério de aceitação depende do grupo: nas lógicas de contatos e bobinas exige-se igualdade bit a bit da imagem de saídas em todas as varreduras; nos temporizadores admite-se erro de até uma varredura mais a resolução da base de tempo; nas grandezas numéricas comparam-se os resultados com o cálculo do oráculo dentro da tolerância do tipo (exata para \textit{S32}, relativa para \textit{F32}).
Os Quadros~\ref{qua:testes-logica} a~\ref{qua:testes-integracao} organizam a cobertura pelo repertório de instruções do Quadro~\ref{qua:grupos-instrucoes}, de modo que cada grupo tenha ao menos um caso associado.
O conjunto completo é conduzido sobre o STM32F767ZI.
A coluna de situação distingue os casos já executados dos apenas especificados: os resultados dos primeiros constam do \ref{ape:pid} e do registro de ensaios do protocolo, e os demais passam a ser registrados nesses quadros à medida que são executados.

\begin{quadro}[htb]
    \centering
    \small\singlespacing\sloppy
    \begin{tabular}{|p{0.8cm}|p{3.2cm}|p{4.8cm}|p{2.8cm}|p{2.3cm}|}
        \hline
        ID & Cobertura & Procedimento & Oráculo & Situação \\
        \hline
        L1 & Contatos e lógica booleana: \texttt{LD}, \texttt{LDN}, \texttt{AND}, \texttt{ANDN}, \texttt{OR}, \texttt{ORN}, \texttt{XOR}, \texttt{XORN}, \texttt{NOT} & Tabela-verdade exaustiva com até três entradas & Tabela-verdade (modelo) & Especificado \\
        \hline
        L2 & Bobinas: \texttt{OUT}, \texttt{SET}, \texttt{RESET}, \texttt{TOGGLE} & Retenção do estado entre varreduras; alternância por borda & Modelo \textit{Stateflow} & Especificado \\
        \hline
        L3 & Bordas: \texttt{LDP}, \texttt{LDF}, \texttt{ANP}, \texttt{ANF}, \texttt{ORP}, \texttt{ORF} & Evento de duração de uma varredura; nível constante não gera evento & Diagrama de tempo & Especificado \\
        \hline
        L4 & Pilhas: \texttt{MPS}, \texttt{MRD}, \texttt{MPP}, \texttt{ANB}, \texttt{ORB} & Ramificações aninhadas até a profundidade máxima da pilha e além dela & Modelo \textit{Stateflow} & Especificado \\
        \hline
        L5 & Lógica de \textit{rung} composta & Selo, intertravamento e partida-parada, combinando L1 a L4 & Modelo \textit{Stateflow} & Especificado \\
        \hline
    \end{tabular}
    \caption{Casos de teste de lógica de contatos e bobinas}
    \label{qua:testes-logica}
    \fonte{Elaborado pelo autor (2026).}
\end{quadro}

\begin{quadro}[htb]
    \centering
    \small\singlespacing
    \begin{tabular}{|p{0.8cm}|p{3.2cm}|p{4.8cm}|p{2.8cm}|p{2.3cm}|}
        \hline
        ID & Cobertura & Procedimento & Oráculo & Situação \\
        \hline
        T1 & \textit{TON} & Tempo programado constante e por registrador; quatro bases de tempo; \textit{EN} interrompido antes do fim; alteração do tempo programado durante a contagem & Diagrama IEC 61131-3 & Especificado \\
        \hline
        T2 & \textit{TOF} & Idem; religamento durante a contagem & Diagrama IEC 61131-3 & Especificado \\
        \hline
        T3 & \textit{TP} & Pulso de duração fixa; reacionamento durante o pulso ignorado & Diagrama IEC 61131-3 & Especificado \\
        \hline
        C1 & \textit{CTU}, \textit{CTD}, \textit{CTUD} & Contagem por borda de subida; saídas nos limites; valor programado constante e por registrador & Norma IEC 61131-3 & Especificado \\
        \hline
        C2 & \textit{Reset} e carga (\textit{load}) & Prioridade entre \textit{reset}, carga e contagem; valor negativo & Norma IEC 61131-3 & Especificado \\
        \hline
    \end{tabular}
    \caption{Casos de teste de temporizadores e contadores}
    \label{qua:testes-temporizacao}
    \fonte{Elaborado pelo autor (2026).}
\end{quadro}

\begin{quadro}[htb]
    \centering
    \small\singlespacing
    \begin{tabular}{|p{0.8cm}|p{3.2cm}|p{4.8cm}|p{2.8cm}|p{2.3cm}|}
        \hline
        ID & Cobertura & Procedimento & Oráculo & Situação \\
        \hline
        N1 & Aritmética \textit{S32}: soma, subtração, multiplicação, divisão, módulo, mínimo, máximo, incremento, decremento, negação e valor absoluto & Valores nominais, zero e limites do inteiro de 32 \textit{bits} & Cálculo em \textit{MATLAB} & Especificado \\
        \hline
        N2 & Aritmética \textit{F32} (mesmas operações) & Valores nominais, muito pequenos e muito grandes; tolerância relativa & Cálculo em \textit{MATLAB} & Especificado \\
        \hline
        N3 & Casos-limite aritméticos & Divisão por zero (resultado nulo), módulo por zero (sem escrita), valor absoluto de \texttt{INT32\_MIN} (saturado), escalonamento com faixa de entrada nula (sem escrita) & Especificação da instrução (\ref{ape:isa}) & Especificado \\
        \hline
        N4 & Comparação \textit{S32} e \textit{F32}: igual, diferente, maior, menor, maior ou igual, menor ou igual & Fronteiras (igualdade, mais e menos uma unidade) & Cálculo em \textit{MATLAB} & Especificado \\
        \hline
        N5 & Conversão e escalonamento: \texttt{ITOF}, \texttt{FTOI}, \texttt{SCALE} & Pontos extremos e intermediários da faixa; operando constante e por registrador & Cálculo em \textit{MATLAB} & Especificado \\
        \hline
    \end{tabular}
    \caption{Casos de teste de blocos numéricos}
    \label{qua:testes-numericos}
    \fonte{Elaborado pelo autor (2026).}
\end{quadro}

\begin{quadro}[htb]
    \centering
    \small\singlespacing
    \begin{tabular}{|p{0.8cm}|p{3.2cm}|p{4.8cm}|p{2.8cm}|p{2.3cm}|}
        \hline
        ID & Cobertura & Procedimento & Oráculo & Situação \\
        \hline
        A1 & Entradas analógicas, saídas PWM e saída analógica & Erro em $N$ pontos da faixa & Instrumento de medição & Especificado \\
        \hline
        B1 & Bloco \textit{RAMP} & \texttt{EN} desligado; borda de subida com valor inicial; subida e descida com taxas distintas; saída no alvo & Modelo \textit{Stateflow} & Especificado \\
        \hline
        B2 & Bloco PID: sete cenários de unidade e regressão & Modos, transferência sem salto, banda morta, \textit{setpoint} dinâmico (\ref{ape:pid}) & Modelo de referência & Executado \\
        \hline
        B3 & Bloco PID em malha fechada (\textit{HIL}) & Cinco arquétipos de planta (\ref{ape:pid}) & Modelo da planta & Executado \\
        \hline
    \end{tabular}
    \caption{Casos de teste de entradas e saídas analógicas e de blocos de controle}
    \label{qua:testes-controle}
    \fonte{Elaborado pelo autor (2026).}
\end{quadro}

\begin{quadro}[htb]
    \centering
    \small\singlespacing
    \begin{tabular}{|p{0.8cm}|p{3.2cm}|p{4.8cm}|p{2.8cm}|p{2.3cm}|}
        \hline
        ID & Cobertura & Procedimento & Oráculo & Situação \\
        \hline
        I1 & Transferência de \textit{bytecode} & Programa recebido idêntico ao enviado; execução após a carga & Comparação byte a byte & Especificado \\
        \hline
        I2 & Robustez do protocolo & Comando inválido, tamanho inconsistente, CRC incorreto e identificador de dispositivo inválido, com o \textit{NACK} correto para cada categoria & Código de \textit{NACK} esperado & Executado \\
        \hline
        I3 & Robustez da máquina virtual & Opcode inválido leva a \textit{trap}; índice fora da faixa não acessa memória externa às regiões & Comportamento especificado & Especificado \\
        \hline
        P1 & Tempo de varredura & Programas de tamanho crescente; média, máximo e variação & Medição & Especificado \\
        \hline
        P2 & Consumo de memória & \textit{Flash} e RAM ocupadas, a partir do relatório de \textit{link} & Medição & Especificado \\
        \hline
        P3 & Custo por grupo de instruções & Tempo de execução por instrução em cada grupo & Medição & Especificado \\
        \hline
    \end{tabular}
    \caption{Casos de teste de integração, robustez e desempenho}
    \label{qua:testes-integracao}
    \fonte{Elaborado pelo autor (2026).}
\end{quadro}


%% ------------------------------------------------------------------------- %%
\section{Síntese}
\label{sec:sintese-implementacao}

Este capítulo apresentou a implementação da solução cujos fundamentos foram estabelecidos no Capítulo~\ref{cap:fundamentacao}: um motor de execução \textit{fetch-decode-execute} com palavra-base de instrução de 32 \textit{bits} e despacho em tempo constante; blocos funcionais com estado segregado em mapas de memória de acesso validado; uma camada de abstração que confina ao contrato de um único módulo toda a dependência do STM32F767ZI; um protocolo com enquadramento COBS e verificação CRC16; modelos de referência em \textit{Simulink}/\textit{Stateflow}; e uma estratégia de validação em quatro níveis, com cobertura de todos os grupos do repertório de instruções.
As escolhas de projeto e suas justificativas estão reunidas no Quadro~\ref{qua:escolhas-projeto}.

\begin{quadro}[htb]
    \centering
    \begin{tabular}{|p{4.7cm}|p{9.7cm}|}
        \hline
        Escolha & Justificativa \\
        \hline
        Palavra-base de 32 \textit{bits} com estrutura auxiliar opcional & Decodificação uniforme da primeira palavra \cite{patterson:2020}; constantes entram como palavras adicionais e parâmetros de blocos por instruções \texttt{ST} e \texttt{STW}, sem ampliar o formato base. \\
        \hline
        Despacho por tabela de opcodes & Tempo de despacho constante e independente da instrução \cite{liu:2000}. \\
        \hline
        Alocação estática de memória & Sem fragmentação nem variação de tempo de alocação \cite{liu:2000}. \\
        \hline
        \textit{Bytecode} somente leitura e estado em mapas & Permite atualizar a lógica sem regravar o \textit{firmware} e conter erros do usuário (\textit{sandbox}). \\
        \hline
        Validação de índice em cada acesso & Um erro de endereçamento não alcança memória fora das regiões alocadas. \\
        \hline
        \textit{Bytecode} validado e pré-processado na carga & Retira configuração estática do laço de varredura. \\
        \hline
        Aritmética \textit{S32} e \textit{F32} nativa & A FPU do alvo executa \textit{F32} em \textit{hardware} \cite{st:ds11532}. \\
        \hline
        Temporização pela base de tempo do sistema & O intervalo independe do tamanho do programa. \\
        \hline
        PID com amostragem fixa & A variação do tempo de varredura não entra nos termos integral e derivativo \cite{astrom:2006}. \\
        \hline
        Derivativo sobre a variável de processo & Evita o pico de saída quando o \textit{setpoint} muda \cite{astrom:2006}. \\
        \hline
        Anti-\textit{windup} por integração condicional & Suspende a integração sob saturação e a retoma na inversão do erro \cite{astrom:2006}. \\
        \hline
        Rastreamento do integrador no modo manual & Transição de manual para automático sem salto \cite{astrom:2006}. \\
        \hline
        Enquadramento COBS e CRC16 & Delimitação sem ambiguidade com sobrecarga limitada e detecção de corrupção \cite{cheshire:1999}. \\
        \hline
        Identificador de dispositivo nos quadros & Rejeição de comandos endereçados a outro dispositivo. \\
        \hline
        \textit{STM32Cube HAL} em vez de acesso direto a registradores & Reduz risco de erro de inicialização; sobrecusto contido na camada de abstração. \\
        \hline
        STM32F767ZI & FPU, memórias acopladas ao núcleo, periféricos e ferramental integrado \cite{st:ds11532,st:um1974}. \\
        \hline
    \end{tabular}

    \caption{Escolhas de projeto e justificativas}
    \label{qua:escolhas-projeto}
    \fonte{Elaborado pelo autor (2026).}
\end{quadro}

% !TeX root = ../main.tex
\chapter{Considerações Finais e Sugestões para Trabalhos Futuros}
\label{cap:finais}

%% ------------------------------------------------------------------------- %%
\section{Conclusões}
\label{sec:conclusoes}

Este trabalho implementou e validou, sobre o microcontrolador STM32F767ZI, uma máquina virtual de lógica \textit{Ladder} cujo motor de execução depende da plataforma apenas pelo contrato de uma camada de abstração de \textit{hardware}.
A solução reúne um conjunto de 76 instruções, com palavra-base de 32 \textit{bits}, que cobre contatos, bobinas, bordas, pilhas lógicas, temporizadores, contadores, aritmética em inteiros e ponto flutuante, escalonamento, entradas e saídas analógicas e os blocos de controle PID e de rampa; memória alocada estaticamente, com validação de índice em cada acesso; e um protocolo serial com enquadramento COBS e verificação CRC16 para a carga do \textit{bytecode} e a depuração.
Cada escolha de projeto foi acompanhada de justificativa (Quadro~\ref{qua:escolhas-projeto}).

A validação seguiu um catálogo de casos de teste organizado por níveis (unidade, integração, sistema e desempenho) e cobrindo todos os grupos do repertório de instruções.
Dos 220 passos de verificação funcional dos ensaios sobre a plataforma-alvo, 219 foram registrados como conformes, o que inclui a igualdade com a tabela-verdade em todas as combinações de entradas da lógica de contatos, a temporização e a contagem dentro das tolerâncias do critério de aceitação, resultados aritméticos exatos em \textit{S32} e dentro da tolerância relativa em \textit{F32}, e o comportamento do bloco PID; e todos os 52 pacotes malformados aplicados ao protocolo foram rejeitados com o código de \textit{NACK} previsto.
O processo de validação também revelou quatro defeitos de implementação e três divergências entre a especificação do conjunto de instruções e o \textit{firmware}, todos corrigidos (Quadro~\ref{qua:defeitos}), o que evidencia o valor de casos de teste que vão além do funcionamento nominal --- programas maiores que um pacote, pacotes malformados e transições de pilha.
Quanto ao desempenho, na configuração de depuração do compilador o tempo de varredura cresce linearmente com o tamanho do programa, a cerca de 37~$\mu$s por \textit{rung} de cinco instruções, e o programa de tamanho máximo suportado leva menos de 10~ms por varredura; o \textit{firmware} ocupa 3,74\% da memória \textit{Flash} e 3,05\% da RAM.

O Quadro~\ref{qua:atendimento-objetivos} relaciona os objetivos específicos ao que foi alcançado.

\begin{quadro}[htb]
    \centering
    \small
    \begin{tabular}{|p{5.4cm}|p{2.2cm}|p{6.6cm}|}
        \hline
        Objetivo específico & Atendimento & Evidência e ressalva \\
        \hline
        Especificar e documentar a \textit{ISA} de 32 \textit{bits} compatível com o \textit{Ladder} da IEC 61131-3 & Atendido & Especificação dos 76 \textit{opcodes} no \ref{ape:isa}, conferida contra o \textit{firmware} e corrigida nos três pontos divergentes. \\
        \hline
        Modelar lógicas de referência em \textit{Simulink}/\textit{Stateflow} & Parcial & Especificação da máquina de modos do bloco PID (Seção~\ref{sec:stateflow-mbd}); as demais lógicas tiveram como referência tabelas-verdade, cálculos independentes e um modelo executável em \textit{Python}. \\
        \hline
        Implementar o motor de execução em linguagem C & Atendido & Seção~\ref{sec:vm-isa}; casos L, T, C, N e B do catálogo. \\
        \hline
        Desenvolver a \textit{HAL} para o \textit{STM32F7} & Atendido & Seção~\ref{sec:hal-implementacao}; uma entrada analógica e as saídas PWM e analógica verificadas sem medição elétrica dos sinais de saída. \\
        \hline
        Validar cruzadamente a \textit{VM} frente ao modelo de referência simulado & Parcial & Comparação com tabelas-verdade, cálculos independentes e modelo de referência em \textit{Python}; a comparação com os modelos \textit{Simulink}/\textit{Stateflow} não foi realizada. \\
        \hline
        Analisar o desempenho (varredura, \textit{RAM} e \textit{Flash}) & Atendido & Seção~\ref{sub:resultados-desempenho}; medido em uma única configuração de compilação (depuração). \\
        \hline
    \end{tabular}

    \caption{Atendimento dos objetivos específicos}
    \label{qua:atendimento-objetivos}
    \fonte{Elaborado pelo autor (2026).}
\end{quadro}

O objetivo geral foi atendido no que concerne à implementação e à validação funcional sobre a plataforma-alvo.
A afirmação de fidelidade à IEC 61131-3 restringe-se ao subconjunto \textit{Ladder} ensaiado, e a de execução determinística se apoia na alocação estática, no despacho em tempo constante e na regularidade do tempo de varredura medido --- a razão entre o máximo e o mínimo ficou entre 1,25 e 1,30 nos programas de 50 a 200 \textit{rungs}, parte dela atribuível à própria variação das entradas durante o acionamento das chaves ---, sem uma análise estatística de variação (\textit{jitter}) da varredura.
As demais limitações dos resultados foram discutidas na Seção~\ref{sub:discussao-resultados}: a ausência da validação cruzada contra os modelos \textit{Simulink}/\textit{Stateflow}; a bancada com uma única entrada analógica e sem medição elétrica de PWM e da saída analógica, que opera de 0 a 3,3~V; a malha fechada com planta simulada não repetida no STM32F767ZI; e a medição de desempenho restrita à configuração de depuração.

%% ------------------------------------------------------------------------- %%
\section{Sugestões para Trabalhos Futuros}
\label{sec:futuros}

As limitações apontadas orientam as seguintes sugestões:

\begin{itemize}
    \item executar os modelos \textit{Simulink}/\textit{Stateflow} da lógica de contatos, da temporização e da contagem e realizar a validação cruzada contra a máquina virtual embarcada, com comparação das imagens de saída varredura a varredura;
    \item repetir no STM32F767ZI os ensaios \textit{HIL} do bloco PID, com medição elétrica da saída analógica e dos sinais PWM por instrumento de aquisição, e acrescentar o estágio de condicionamento para a faixa de 0 a 10~V usual em CLPs;
    \item medir o tempo de varredura com otimização do compilador e com a memória de instruções em cache (acelerador de \textit{Flash} e \textit{cache} do núcleo) habilitados, e investigar o comportamento do programa de 10 \textit{rungs}, que não segue o ajuste linear dos demais;
    \item caracterizar a variação do tempo de varredura (\textit{jitter}) com entradas estáveis e sob carga de comunicação, para sustentar quantitativamente a afirmação de determinismo;
    \item ampliar a bancada às três entradas analógicas e a mais canais digitais, e executar ensaios de longa duração e de injeção de falhas na comunicação;
    \item avaliar o desempenho e a aderência do conjunto de instruções em outras plataformas, por meio de novas implementações do contrato da camada de abstração.
\end{itemize}


% ---------------------------------------------------------------------------- %
% Apêndices e/ou Anexos

\appendix
\renewcommand{\appendixname}{Apêndice} % Para trocar Apendice por Anexo (Se necessário)
\setcounter{chapter}{0} % ressetar a contagem

% Numeração dos apêndices (NBR 14724): "APÊNDICE A --- TÍTULO"; seções, tabelas
% e figuras internas numeradas como A.1, A.2 etc.
\renewcommand{\thechapter}{\appendixname~\Alph{chapter}}
\renewcommand{\thesection}{\Alph{chapter}.\arabic{section}}
\renewcommand{\thetable}{\Alph{chapter}.\arabic{table}}
\renewcommand{\thefigure}{\Alph{chapter}.\arabic{figure}}
\renewcommand{\thequadro}{\Alph{chapter}.\arabic{quadro}}
\titleformat{\chapter}[hang]{\normalfont\bfseries}{\MakeUppercase{\thechapter}\ ---}{0.5em}{\MakeUppercase}

% !TeX root = ../main.tex
\addtocontents{toc}{\protect\let\protect\numberline\protect\apendiceNumberline}
\chapter{Especificação do Conjunto de Instruções}
\label{ape:isa}
\addtocontents{toc}{\protect\let\protect\numberline\protect\numberlineorig}

Este apêndice detalha o conjunto de instruções da máquina virtual apresentado na Seção~\ref{sec:vm-isa}: convenções de codificação, tipos de operando, sub-tipos dos blocos funcionais e a descrição de cada um dos 76 opcodes.

% ---------------------------------------------------------------------------- %
\section{Convenções de codificação}
\label{sec:ape-isa-convencoes}

Toda instrução ocupa uma palavra de 4 \textit{bytes}, alinhada em palavra; instruções com operando imediato de 32 \textit{bits} ocupam uma palavra adicional imediatamente após a primeira, com o valor em ordem \textit{big-endian}.
O \textit{byte} 0 contém o opcode nos 7 \textit{bits} menos significativos; o \textit{bit} 7 é reservado à depuração e mascarado antes do despacho.
No formato genérico de operando (Figura~\ref{fig:formato-instrucao}), o \textit{byte} 1 contém o tipo do operando (4 \textit{bits}) e a parte alta do índice (4 \textit{bits}), e o \textit{byte} 2 contém a parte baixa do índice (4 \textit{bits}) e o sub-tipo (4 \textit{bits}), o que permite endereçar até 256 elementos por tipo; o \textit{byte} 3 é reservado, exceto nas instruções aritméticas, de conversão e de escalonamento, nas quais contém o índice do registrador de destino, e em \texttt{STW}, na qual contém o índice do bloco de destino.
As instruções \texttt{PID}, \texttt{RAMP} e \texttt{PWM\_OUT} utilizam apenas o \textit{byte} 1 (índice da instância e base de tempo).
Um opcode superior a \texttt{0x4B}, exceto \texttt{0x7F}, é inválido e provoca o estado de erro (\textit{trap}).

A Tabela~\ref{tab:ape-tipos} lista os tipos de operando, e a Tabela~\ref{tab:ape-subtipos} os sub-tipos que selecionam o membro de um bloco funcional.

\begin{table}[htb]
    \centering
    \caption{Tipos de operando}
    \label{tab:ape-tipos}
    \begin{tabular}{|c|l|p{7.6cm}|}
        \hline
        Código & Tipo & Descrição \\
        \hline
        \texttt{0x01} & \texttt{\%I} & Entrada digital física \\
        \texttt{0x02} & \texttt{\%Q} & Saída digital física \\
        \texttt{0x03} & \texttt{\%M} & Memória auxiliar (\textit{bit}) \\
        \texttt{0x04} & \texttt{\%T} & Temporizador \\
        \texttt{0x05} & \texttt{\%C} & Contador \\
        \texttt{0x06} & \texttt{DR} & Registrador de dados de 32 \textit{bits} \\
        \texttt{0x07} & \texttt{CONST} & Constante imediata (próxima palavra) \\
        \texttt{0x08} & \texttt{\%AN} & Entrada analógica (somente \texttt{MOV} e \texttt{SCALE}) \\
        \texttt{0x09} & \texttt{\%P} & Saída PWM (canais 1 a 3) \\
        \texttt{0x0A} & \texttt{\%AO} & Saída analógica pura \\
        \texttt{0x0B} & \texttt{\%PID} & Instância do bloco PID \\
        \texttt{0x0C} & \texttt{\%RAMP} & Instância do bloco RAMP \\
        \hline
    \end{tabular}

    {\footnotesize Fonte: Elaborado pelo autor (2026).\par}
\end{table}

\begin{table}[htb]
    \centering
    \caption{Sub-tipos de operando dos blocos funcionais}
    \label{tab:ape-subtipos}
    \begin{tabular}{|l|p{3.1cm}|p{7.3cm}|}
        \hline
        Bloco & Sub-tipo & Valores \\
        \hline
        Temporizador & \textit{Bit} / palavra & \texttt{EN}=0, \texttt{Q}=1 / \texttt{ET}=2 (tempo decorrido), \texttt{PT}=3 (tempo programado) \\
        \hline
        Contador & \textit{Bit} / palavra & \texttt{EN}=0, \texttt{CU}=1, \texttt{CD}=2, \texttt{RST}=3, \texttt{LOAD}=4, \texttt{QU}=5, \texttt{QD}=6 / \texttt{CV}=7 (valor atual), \texttt{PV}=8 (valor programado) \\
        \hline
        PID & \textit{Bit} (entradas) & \texttt{EN}=0, \texttt{RESET}=1, \texttt{MANUAL}=2 \\
        & \textit{Bit} (saídas) & \texttt{Q}=3, \texttt{AT\_MAX}=4, \texttt{AT\_MIN}=5, \texttt{MAN\_O}=6 \\
        & Palavra F32 (saídas) & \texttt{CV}=0, \texttt{ERROR}=1, \texttt{INTEGRAL}=2 \\
        \hline
        RAMP & \textit{Bit} & \texttt{EN}=0 (entrada), \texttt{Q}=1 (saída: \texttt{OUT} igual a \texttt{TARGET}) \\
        & Palavra F32 & \texttt{OUT}=0, \texttt{TARGET}=1, \texttt{INIT}=2, \texttt{ASCEND}=3, \texttt{DESCEND}=4 \\
        \hline
        PWM / saída analógica & Palavra & \texttt{DUTY} (0 a 10000, resolução de 0,01\%) e \texttt{FREQ} (10 a 1000 Hz) \\
        \hline
    \end{tabular}

    {\footnotesize Fonte: Elaborado pelo autor (2026).\par}
\end{table}

Os parâmetros de entrada de \texttt{PWM}, \texttt{PID} e \texttt{RAMP} são escritos pela instrução \texttt{STW}, que grava em um membro de bloco o valor de uma fonte (registrador ou constante \textit{IEEE-754}) de forma incondicional a cada varredura; a Tabela~\ref{tab:ape-stw} lista os destinos.

\begin{table}[htb]
    \centering
    \caption{Destinos da instrução \texttt{STW}}
    \label{tab:ape-stw}
    \begin{tabular}{|c|p{5.5cm}|c|p{5.5cm}|}
        \hline
        Código & Destino & Código & Destino \\
        \hline
        \texttt{0x00} & PWM: \textit{duty} (0 a 10000) & \texttt{0x09} & PID: limite máximo de saída \\
        \texttt{0x01} & PWM: frequência (10 a 1000 Hz) & \texttt{0x0A} & PID: valor manual de \texttt{CV} \\
        \texttt{0x02} & PID: \textit{setpoint} & \texttt{0x0B} & PID: tempo de amostragem \\
        \texttt{0x03} & PID: variável de processo & \texttt{0x0C} & RAMP: valor alvo \\
        \texttt{0x04} & PID: ganho proporcional & \texttt{0x0D} & RAMP: taxa máxima de subida \\
        \texttt{0x05} & PID: ganho integral (1/s) & \texttt{0x0E} & RAMP: taxa máxima de descida \\
        \texttt{0x06} & PID: tempo derivativo (s) & \texttt{0x0F} & RAMP: múltiplo da base de tempo \\
        \texttt{0x07} & PID: banda morta & \texttt{0x10} & RAMP: valor inicial \\
        \texttt{0x08} & PID: limite mínimo de saída & & \\
        \hline
    \end{tabular}

    {\footnotesize Fonte: Elaborado pelo autor (2026).\par}
\end{table}

% ---------------------------------------------------------------------------- %
\section{Instruções de controle de fluxo e pilhas}
\label{sec:ape-isa-fluxo}

A máquina mantém um acumulador booleano e uma pilha lógica de 16 posições, para ramificações, e uma pilha de blocos de 16 posições, para a associação de trechos de \textit{rung}.

\begingroup\singlespacing
\begin{longtable}{|l|c|p{9.4cm}|}
    \caption{Instruções de controle de fluxo e pilhas}
    \label{tab:ape-fluxo}\\
    \hline
    Mnemônico & Opcode & Operação \\
    \hline
    \endfirsthead
    \hline
    Mnemônico & Opcode & Operação \\
    \hline
    \endhead
    \hline
    \endfoot
    \texttt{NOP} & \texttt{0x00} & Nenhuma operação; preserva o alinhamento do programa. \\
    \texttt{SOR} & \texttt{0x01} & Início de \textit{rung}: zera as pilhas e inicializa o acumulador em verdadeiro; o \textit{byte} 1 carrega o número da \textit{rung}, usado em depuração. \\
    \texttt{MPS} & \texttt{0x12} & Empilha o acumulador na pilha lógica; provoca \textit{trap} se a pilha estiver cheia. \\
    \texttt{MRD} & \texttt{0x13} & Copia o topo da pilha lógica para o acumulador, sem desempilhar. \\
    \texttt{MPP} & \texttt{0x14} & Desempilha o topo da pilha lógica para o acumulador. \\
    \texttt{ANB} & \texttt{0x15} & \textit{E} lógico entre o valor desempilhado da pilha de blocos e o acumulador (circuitos em série). \\
    \texttt{ORB} & \texttt{0x16} & \textit{OU} lógico entre o valor desempilhado da pilha de blocos e o acumulador (circuitos em paralelo). \\
    \texttt{EOS} & \texttt{0x7F} & Fim de varredura: encerra o laço de execução. \\
\end{longtable}
\endgroup

% ---------------------------------------------------------------------------- %
\section{Contatos, lógica booleana e bobinas}
\label{sec:ape-isa-bit}

As instruções de \textit{bit} aceitam como operando os tipos \texttt{\%I}, \texttt{\%Q}, \texttt{\%M}, e os membros booleanos de temporizadores, contadores, blocos PID e RAMP.
As instruções sensíveis a borda comparam o valor atual do operando com o da varredura anterior, mantido em um mapa de últimos valores.

\begingroup\singlespacing
\begin{longtable}{|l|c|p{9.4cm}|}
    \caption{Contatos, lógica booleana e bobinas}
    \label{tab:ape-bit}\\
    \hline
    Mnemônico & Opcode & Operação \\
    \hline
    \endfirsthead
    \hline
    Mnemônico & Opcode & Operação \\
    \hline
    \endhead
    \hline
    \endfoot
    \texttt{ST} & \texttt{0x02} & Armazena o acumulador diretamente em um membro de bloco funcional (por exemplo, \texttt{EN} de um temporizador). \\
    \texttt{LD} & \texttt{0x03} & Carrega o operando no acumulador. \\
    \texttt{LDN} & \texttt{0x04} & Carrega o complemento do operando. \\
    \texttt{LDP} & \texttt{0x05} & Carrega verdadeiro se o operando fez transição de 0 para 1 desde a varredura anterior. \\
    \texttt{LDF} & \texttt{0x06} & Carrega verdadeiro se o operando fez transição de 1 para 0. \\
    \texttt{AND} & \texttt{0x07} & \textit{E} entre o acumulador e o operando. \\
    \texttt{ANDN} & \texttt{0x08} & \textit{E} entre o acumulador e o complemento do operando. \\
    \texttt{ANP} & \texttt{0x09} & \textit{E} entre o acumulador e a borda de subida do operando. \\
    \texttt{ANF} & \texttt{0x0A} & \textit{E} entre o acumulador e a borda de descida do operando. \\
    \texttt{OR} & \texttt{0x0B} & \textit{OU} entre o acumulador e o operando. \\
    \texttt{ORN} & \texttt{0x0C} & \textit{OU} entre o acumulador e o complemento do operando. \\
    \texttt{ORP} & \texttt{0x0D} & \textit{OU} entre o acumulador e a borda de subida do operando. \\
    \texttt{ORF} & \texttt{0x0E} & \textit{OU} entre o acumulador e a borda de descida do operando. \\
    \texttt{NOT} & \texttt{0x0F} & Inverte o acumulador. \\
    \texttt{XOR} & \texttt{0x10} & \textit{OU} exclusivo entre o acumulador e o operando. \\
    \texttt{XORN} & \texttt{0x11} & \textit{OU} exclusivo entre o acumulador e o complemento do operando. \\
    \texttt{OUT} & \texttt{0x17} & Escreve o acumulador em uma saída ou memória auxiliar. \\
    \texttt{SET} & \texttt{0x18} & Se o acumulador for verdadeiro, escreve 1 no operando; caso contrário, ignora. Para saídas, o valor é retido entre varreduras. \\
    \texttt{RESET} & \texttt{0x19} & Se o acumulador for verdadeiro, escreve 0 no operando; caso contrário, ignora. \\
    \texttt{TOGGLE} & \texttt{0x1A} & Inverte o operando a cada transição do acumulador de 0 para 1. \\
\end{longtable}
\endgroup

% ---------------------------------------------------------------------------- %
\section{Temporizadores e contadores}
\label{sec:ape-isa-tc}

\begingroup\singlespacing
\begin{longtable}{|l|c|p{9.4cm}|}
    \caption{Instruções de temporização e contagem}
    \label{tab:ape-tc}\\
    \hline
    Mnemônico & Opcode & Operação \\
    \hline
    \endfirsthead
    \hline
    Mnemônico & Opcode & Operação \\
    \hline
    \endhead
    \hline
    \endfoot
    \texttt{TIM} & \texttt{0x1B} & Executa uma instância de temporizador (32 instâncias). Campos: tipo (\textit{TON}, \textit{TOF} ou \textit{TP}), índice da instância, origem do tempo programado (constante imediata ou registrador) e base de tempo (10 ms, 100 ms, 1 s ou 1 min). O bit \texttt{EN} é armazenado previamente por \texttt{ST}. \textit{TON}: \texttt{Q} verdadeiro após o tempo programado decorrido desde a borda de subida de \texttt{EN}. \textit{TOF}: \texttt{Q} verdadeiro enquanto \texttt{EN} verdadeiro e por mais o tempo programado após a borda de descida. \textit{TP}: pulso de duração programada a partir da borda de subida, não retentável. \\
    \texttt{COUNT} & \texttt{0x1C} & Executa uma instância de contador (32 instâncias). Campos: tipo (\textit{CTU}, \textit{CTD} ou \textit{CTUD}), índice da instância e origem do valor programado (constante imediata ou registrador). Entradas \texttt{EN}, \texttt{CU}, \texttt{CD}, \texttt{RST} e \texttt{LOAD} armazenadas por \texttt{ST}; saídas \texttt{QU} (valor atual maior ou igual ao programado) e \texttt{QD} (valor atual não positivo); valor atual de 32 \textit{bits} com sinal. \\
\end{longtable}
\endgroup

% ---------------------------------------------------------------------------- %
\section{Aritmética, comparação, conversão e escalonamento}
\label{sec:ape-isa-num}

As operações numéricas existem em duas versões, \textit{S32} (inteiro de 32 \textit{bits} com sinal) e \textit{F32} (ponto flutuante de precisão simples); os dois opcodes de cada linha da Tabela~\ref{tab:ape-aritmetica} são apresentados nessa ordem.
As operações binárias seguem o padrão \texttt{LDW} (carrega o primeiro operando no acumulador de palavra) seguido da operação com um segundo operando e gravação em um registrador de destino.
Fontes aceitas em \textit{S32}: registrador, constante, temporizador e contador; em \textit{F32}: registrador, constante, blocos PID e RAMP (para usar um inteiro em contexto \textit{F32} aplica-se antes \texttt{ITOF}).

\begingroup\singlespacing
\begin{longtable}{|l|c|p{9.4cm}|}
    \caption{Aritmética e movimentação}
    \label{tab:ape-aritmetica}\\
    \hline
    Mnemônico & Opcodes & Operação \\
    \hline
    \endfirsthead
    \hline
    Mnemônico & Opcodes & Operação \\
    \hline
    \endhead
    \hline
    \endfoot
    \texttt{LDW} & \texttt{0x1D}/\texttt{0x1E} & Carrega o operando no acumulador de palavra. \\
    \texttt{ADD} & \texttt{0x1F}/\texttt{0x20} & acumulador $+$ operando $\rightarrow$ destino. \\
    \texttt{SUB} & \texttt{0x21}/\texttt{0x22} & acumulador $-$ operando $\rightarrow$ destino. \\
    \texttt{MUL} & \texttt{0x23}/\texttt{0x24} & acumulador $\times$ operando $\rightarrow$ destino. \\
    \texttt{DIV} & \texttt{0x25}/\texttt{0x26} & acumulador $\div$ operando $\rightarrow$ destino; o resultado é zero se o divisor for zero. \\
    \texttt{MOV} & \texttt{0x27}/\texttt{0x28} & operando $\rightarrow$ destino (\textit{S32} aceita também entrada analógica). \\
    \texttt{INC} & \texttt{0x29}/\texttt{0x2A} & operando $+ 1 \rightarrow$ destino. \\
    \texttt{DEC} & \texttt{0x2B}/\texttt{0x2C} & operando $- 1 \rightarrow$ destino. \\
    \texttt{NEG} & \texttt{0x2D}/\texttt{0x2E} & $-$operando $\rightarrow$ destino. \\
    \texttt{ABS} & \texttt{0x2F}/\texttt{0x30} & $\lvert$operando$\rvert \rightarrow$ destino; em \textit{S32}, o menor inteiro é saturado no maior. \\
    \texttt{MOD} & \texttt{0x31}/\texttt{0x32} & Resto da divisão (\texttt{fmodf} em \textit{F32}) $\rightarrow$ destino; sem escrita se o divisor for zero. \\
    \texttt{MIN} & \texttt{0x33}/\texttt{0x34} & Mínimo entre o acumulador e o operando $\rightarrow$ destino. \\
    \texttt{MAX} & \texttt{0x35}/\texttt{0x36} & Máximo entre o acumulador e o operando $\rightarrow$ destino. \\
\end{longtable}
\endgroup

As instruções de comparação avaliam o acumulador de palavra contra um operando: se a relação for verdadeira, o acumulador booleano permanece inalterado; caso contrário, é zerado, o que permite compor comparações com contatos em série (Tabela~\ref{tab:ape-comparacao}).

\begingroup\singlespacing
\begin{longtable}{|l|c|p{9.4cm}|}
    \caption{Comparação, conversão e escalonamento}
    \label{tab:ape-comparacao}\\
    \hline
    Mnemônico & Opcodes & Operação \\
    \hline
    \endfirsthead
    \hline
    Mnemônico & Opcodes & Operação \\
    \hline
    \endhead
    \hline
    \endfoot
    \texttt{EQU} & \texttt{0x37}/\texttt{0x38} & Igual. \\
    \texttt{NEQ} & \texttt{0x39}/\texttt{0x3A} & Diferente. \\
    \texttt{GRT} & \texttt{0x3B}/\texttt{0x3C} & Maior que. \\
    \texttt{LES} & \texttt{0x3D}/\texttt{0x3E} & Menor que. \\
    \texttt{GEQ} & \texttt{0x3F}/\texttt{0x40} & Maior ou igual. \\
    \texttt{LEQ} & \texttt{0x41}/\texttt{0x42} & Menor ou igual. \\
    \texttt{ITOF} & \texttt{0x43} & Converte inteiro em real: $(\text{f32})\,\text{operando} \rightarrow$ destino. \\
    \texttt{FTOI} & \texttt{0x44} & Converte real em inteiro: $(\text{s32})\,\text{operando} \rightarrow$ destino. \\
    \texttt{SCALE} & \texttt{0x45}/\texttt{0x46} & Escalonamento linear $out_{min} + (X - in_{min})\,(out_{max} - out_{min}) / (in_{max} - in_{min}) \rightarrow$ destino. $X$ pode ser registrador, constante ou entrada analógica; os quatro limites são constantes imediatas de 32 \textit{bits} (instrução de 20 \textit{bytes}); sem escrita se $in_{min} = in_{max}$. \\
\end{longtable}
\endgroup

% ---------------------------------------------------------------------------- %
\section{Saídas analógicas, escrita em bloco e blocos de controle}
\label{sec:ape-isa-blocos}

\begingroup\singlespacing
\begin{longtable}{|l|c|p{9.4cm}|}
    \caption{Saídas analógicas, escrita em bloco e blocos de controle}
    \label{tab:ape-blocos}\\
    \hline
    Mnemônico & Opcode & Operação \\
    \hline
    \endfirsthead
    \hline
    Mnemônico & Opcode & Operação \\
    \hline
    \endhead
    \hline
    \endfoot
    \texttt{ANALOG\_OUT} & \texttt{0x47} & Escreve uma tensão (registrador ou constante \textit{IEEE-754}, com conversão automática de inteiro) na saída analógica pura, limitada a 0 a 10 V. Executa apenas com o acumulador verdadeiro. \\
    \texttt{STW} & \texttt{0x48} & Grava o valor de uma fonte (registrador ou constante) em um membro de bloco funcional (Tabela~\ref{tab:ape-stw}). Executa a cada varredura, independentemente do acumulador. \\
    \texttt{PWM\_OUT} & \texttt{0x49} & Aplica ao canal PWM indicado o \textit{duty} (0 a 100\%, limitado) e a frequência previamente gravados por \texttt{STW}. Executa apenas com o acumulador verdadeiro. \\
    \texttt{PID} & \texttt{0x4A} & Executa uma instância do controlador PID posicional (16 instâncias). O período de amostragem é fixo, dado pelo produto do parâmetro de amostragem pela base de tempo da instrução, e o bloco recalcula apenas quando esse período decorre (Seção~\ref{sub:blocos-funcionais}); prioridade de modos \texttt{EN} $>$ \texttt{RESET} $>$ \texttt{MANUAL} $>$ automático. \\
    \texttt{RAMP} & \texttt{0x4B} & Executa uma instância do limitador de taxa de variação (16 instâncias): a saída persegue o valor alvo com taxa máxima de subida ou de descida por unidade de base de tempo; com \texttt{EN} desligado a saída é zero e, na borda de subida de \texttt{EN}, assume o valor inicial. \\
\end{longtable}
\endgroup

% !TeX root = ../main.tex
\addtocontents{toc}{\protect\let\protect\numberline\protect\apendiceNumberline}
\chapter{Ensaios do Bloco de Controle PID}
\label{ape:pid}
\addtocontents{toc}{\protect\let\protect\numberline\protect\numberlineorig}

Este apêndice registra os ensaios já realizados sobre o bloco PID descrito na Seção~\ref{sub:blocos-funcionais} e cuja máquina de modos é especificada na Seção~\ref{sec:stateflow-mbd}: testes de unidade e regressão e ensaios em malha fechada com \textit{Hardware-in-the-Loop} (\textit{HIL}).
Ambos foram executados na plataforma de desenvolvimento do firmware.
Como o motor de execução da máquina virtual não acessa periféricos fora do contrato da camada de abstração (Seção~\ref{sec:hal-implementacao}), o comportamento validado independe do microcontrolador; a repetição dos ensaios no STM32F767ZI será registrada neste apêndice quando executada.

% ---------------------------------------------------------------------------- %
\section{Testes de unidade e regressão}
\label{sec:ape-pid-unidade}

Os testes de unidade e regressão cobrem sete cenários, identificados de A1 a A7 (caso B2 da Tabela~\ref{tab:catalogo-testes}), cada um definido por objetivo, procedimento numerado e critério de aceite, e executados diretamente na bancada, sem planta simulada.
A Tabela~\ref{tab:ape-unidade-pid} resume os cenários.

\begin{table}[htb]
    \centering
    \caption{Testes de unidade e regressão do bloco PID}
    \label{tab:ape-unidade-pid}
    \begin{tabular}{|l|p{4.6cm}|p{7.6cm}|}
        \hline
        ID & Cenário & Critério de aceite \\
        \hline
        A1 & \texttt{EN} armazenado por \texttt{ST} & \texttt{EN} e \texttt{Q} caem na varredura seguinte ao desligamento; \texttt{CV} vai ao limite mínimo; \texttt{INTEGRAL} é preservado. \\
        A2 & \texttt{CV} sem truncamento de 4 \textit{bits} & A saída analógica é sempre igual a \texttt{CV}, sem colisão de índice de sub-tipo. \\
        A3 & \texttt{RESET} & Integrador zerado e \texttt{CV} igual ao limite mínimo durante o \texttt{RESET}; retomada normal ao soltar. \\
        A4 & Transferência sem salto (manual $\leftrightarrow$ automático) & \texttt{CV} não apresenta salto abrupto em nenhuma das duas transições. \\
        A5 & Banda morta & \texttt{CV} permanece constante enquanto $\lvert SP - PV \rvert$ está dentro da banda. \\
        A6 & \textit{Setpoint} dinâmico & \texttt{PV} converge ao novo \texttt{SP} após cada variação, sem oscilação sustentada. \\
        A7 & Ausência de colisão entre campos do PID & Cada campo do protocolo de depuração retorna um valor distinto e plausível. \\
        \hline
    \end{tabular}

    {\footnotesize Fonte: Elaborado pelo autor (2026).\par}
\end{table}

Os cenários A1 e A2 correspondem à correção de dois defeitos identificados durante o desenvolvimento do bloco.
No primeiro, o bit \texttt{EN} era lido diretamente do acumulador lógico no instante da instrução \texttt{PID}, herdando o resultado de qualquer lógica anterior que reutilizasse o acumulador; a correção adotou o padrão já empregado por temporizadores e contadores, no qual o bit é armazenado por uma instrução \texttt{ST} dedicada, em \textit{rung} própria, e lido do mapa de memória do bloco.
No segundo, o campo de sub-tipo, de apenas 4 \textit{bits} (16 valores), era indexado por um enumerador de 20 membros, de modo que o índice de \texttt{CV} colidia, por aritmética de módulo, com o de \texttt{RESET}; a correção dividiu o enumerador em três --- um para instruções de \textit{bit} (7 valores), um para instruções de palavra (13 valores) e um exclusivo do protocolo de depuração (20 valores, sem a limitação de 4 \textit{bits}) ---, o que é possível porque a família da instrução já é conhecida na decodificação.

% ---------------------------------------------------------------------------- %
\section{Ensaios em malha fechada (Hardware-in-the-Loop)}
\label{sec:ape-pid-hil}

Os ensaios em malha fechada validam o bloco PID com o firmware em execução no microcontrolador, ligado por um instrumento de aquisição analógica (\textit{Analog Discovery 3}) a uma planta simulada em tempo real no \textit{MATLAB}.
O sinal de controle (\texttt{CV}) sai por uma saída analógica real e retorna, após ser processado pelo modelo da planta, por uma entrada analógica real; a física do processo é o único elemento virtual, o que caracteriza o ensaio como \textit{HIL} \cite{isermann:1999}.
O deslocamento de 0,3 V do conversor digital-analógico do instrumento é compensado no \textit{script} de ensaio.
A Figura~\ref{fig:malha-hil} ilustra a malha.

\begin{figure}[htb]
    \centering
    \begin{tikzpicture}[x=1cm,y=1cm]
        \node[bloco, minimum width=3.0cm, minimum height=14mm, font=\footnotesize] (fw) at (0,0) {Máquina virtual\\(bloco PID no\\\textit{firmware})};
        \node[bloco, minimum width=3.0cm, minimum height=14mm, font=\footnotesize] (ad) at (5.2,0) {\textit{Analog}\\\textit{Discovery 3}\\(saída e entrada)};
        \node[bloco, minimum width=3.0cm, minimum height=14mm, font=\footnotesize] (pl) at (10.4,0) {\textit{MATLAB}:\\modelo da\\planta};
        \draw[seta] ([yshift=3mm]fw.east) -- node[rotulo, above=1pt, fill=none] {\texttt{CV}} ([yshift=3mm]ad.west);
        \draw[seta] ([yshift=3mm]ad.east) -- node[rotulo, above=1pt, fill=none] {\texttt{CV}} ([yshift=3mm]pl.west);
        \draw[seta] ([yshift=-3mm]pl.west) -- node[rotulo, below=1pt, fill=none] {\texttt{PV}} ([yshift=-3mm]ad.east);
        \draw[seta] ([yshift=-3mm]ad.west) -- node[rotulo, below=1pt, fill=none] {\texttt{PV}} ([yshift=-3mm]fw.east);
    \end{tikzpicture}
    \caption{Malha de ensaio \textit{Hardware-in-the-Loop}}
    \label{fig:malha-hil}

    {\footnotesize Fonte: Elaborado pelo autor (2026).\par}
\end{figure}

A malha é fechada pelo bloco PID do firmware, e não por um bloco de controle genérico do \textit{Simulink}, para que o algoritmo ensaiado seja exatamente o implementado --- com amostragem fixa, derivativo sobre a variável medida, anti-\textit{windup} por integração condicional e limitação de saída --- e não uma aproximação dele.

O mesmo controlador, com ganhos fixos $K_p = 0{,}4$, $K_i = 0{,}05$ e $K_d = 0$, foi ensaiado contra cinco arquétipos clássicos de processo, escolhidos por cobrirem, em conjunto, as dinâmicas usuais da prática de controle: primeira ordem rápida, primeira ordem lenta (para expor o \textit{windup} integral), primeira ordem com atraso de transporte (o modelo mais comum de processo industrial), integrador puro (processo sem autorregulação) e segunda ordem subamortecida (processos mecânicos ressonantes).
A escolha de $K_d = 0$ é deliberada: o termo derivativo, calculado sobre a variável medida, amplifica diretamente o ruído do conversor analógico-digital \cite{astrom:2006}.
O critério de acomodação é o primeiro instante, após o acionamento do controlador, a partir do qual $\lvert PV - SP \rvert \leq 2$ e a condição não é mais violada até o fim da janela de observação.
A Tabela~\ref{tab:resultados-hil} apresenta os resultados, calculados automaticamente a partir dos registros dos ensaios.

\begin{table}[htb]
    \centering
    \caption{Resultados dos ensaios HIL para os cinco arquétipos de planta}
    \label{tab:resultados-hil}
    \begin{tabular}{|p{6.4cm}|c|p{4.2cm}|}
        \hline
        Arquétipo de planta & \textit{Overshoot} máximo & Tempo de acomodação \\
        \hline
        1 --- Primeira ordem rápida & 0,3\% & 4,0 s \\
        2 --- Primeira ordem lenta & 26,7\% & 76,1 s \\
        3 --- Atraso de transporte (ganhos originais) & 47,2\% & 80,1 s (não convergiu na janela) \\
        4 --- Integrador puro & 6,2\% & 34,0 s \\
        5 --- Segunda ordem subamortecida & 42,2\% & 28,9 s \\
        \hline
    \end{tabular}

    {\footnotesize Fonte: Elaborado pelo autor (2026).\par}
\end{table}

O arquétipo de atraso de transporte (planta 3) apresentou o pior caso da bateria, com \textit{overshoot} de 47,2\% e sinal de ciclo-limite, com a saída de controle saturando repetidamente em zero.
O resultado é esperado teoricamente, já que o atraso de transporte reduz a margem de fase mais severamente do que a lentidão isolada (planta 2) ou a ausência de autorregulação (planta 4), mas a sua magnitude só foi conhecida pelo ensaio.
Aplicou-se então um reajuste de ganhos a essa planta, com base na heurística clássica $T_i \approx \tau + 2L$ para processos com atraso de transporte: $K_p = 0{,}4$ e $K_i = 0{,}05$ foram alterados para $K_p = 0{,}3$ e $K_i = 0{,}02$, o que reduziu o \textit{overshoot} de 47,2\% para 0,10\% e o tempo de acomodação de 80,1~s (sem convergência plena na janela) para 16,0~s.
O caso evidencia que a bateria de cinco arquétipos expõe, de forma quantificada, uma fragilidade de sintonia que a validação em malha aberta ou com uma única planta não revelaria: o mesmo controlador, com os mesmos ganhos, comporta-se de forma satisfatória em quatro das cinco plantas e é criticamente subamortecido na quinta, o que confirma que a sintonia de um PID depende do processo, e não apenas do algoritmo.

% !TeX root = ../main.tex
\addtocontents{toc}{\protect\let\protect\numberline\protect\apendiceNumberline}
\chapter{Ensaio de Robustez do Protocolo de Comunicação}
\label{ape:protocolo}
\addtocontents{toc}{\protect\let\protect\numberline\protect\numberlineorig}

Este apêndice registra o ensaio de robustez do protocolo de comunicação (caso I2 do Quadro~\ref{qua:testes-integracao}), executado sobre o STM32F767ZI em 7 de outubro de 2026, por meio de um \textit{script} no \textit{host} conectado à porta serial virtual do ST-LINK.
O \textit{script} estabelece a conexão, leva o dispositivo a cada um dos quatro modos do protocolo --- conexão, ocioso, transferência e depuração (Seção~\ref{sec:protocolo}) ---, envia em cada um pacotes malformados e compara o código de \textit{NACK} recebido com o previsto para a categoria de erro.
As categorias são: comando inexistente no modo; tamanho do pacote inconsistente com o do comando; CRC16 incorreto; identificador de dispositivo diferente do do alvo; e, para os comandos que a têm, violação de faixa de um campo (carga útil de \texttt{DATA\_PACKET} acima de 256 \textit{bytes}, nível de \textit{debounce} fora de 1 a 10 e 17ª variável na lista de monitoração, que comporta 16).
Cada categoria de cada comando tem um código de \textit{NACK} próprio, de modo que a resposta identifica sem ambiguidade a condição rejeitada.

Os 52 casos foram rejeitados com o código esperado, e os Quadros~\ref{qua:ape-i2-conexao} a~\ref{qua:ape-i2-depuracao} apresentam os códigos obtidos por modo.
Uma execução anterior, com uma versão do \textit{firmware} cuja máquina de estados de conexão tratava de modo incorreto os pacotes inválidos, falhou nos sete casos do modo de conexão (Quadro~\ref{qua:defeitos}); a execução apresentada é a posterior à correção.

% ---------------------------------------------------------------------------- %
\section{Modos de conexão e ocioso}
\label{sec:ape-i2-conexao}

\begin{quadro}[htb]
    \centering
    \small
    \begin{tabular}{|l|c|c|c|c|}
        \hline
        Comando & Tamanho & CRC & Identificador & Outras condições \\
        \hline
        Comando inexistente (conexão) & --- & --- & --- & 0x60 \\
        \hline
        \texttt{START\_CONNECTION} & 0x61 & 0x62 & 0x63 & --- \\
        \hline
        \texttt{VALIDATION} & 0x64 & 0x65 & 0x66 & --- \\
        \hline
    \end{tabular}

    \caption{Códigos de \textit{NACK} obtidos no modo de conexão (7 casos)}
    \label{qua:ape-i2-conexao}
    \fonte{Elaborado pelo autor (2026).}
\end{quadro}

\begin{quadro}[htb]
    \centering
    \small
    \begin{tabular}{|l|c|c|c|c|}
        \hline
        Comando & Tamanho & CRC & Identificador & Outras condições \\
        \hline
        Comando inexistente (ocioso) & --- & --- & --- & 0x00 \\
        \hline
        \texttt{START\_TRANSFER} & 0x01 & 0x03 & 0x05 & --- \\
        \hline
        \texttt{START\_DEBUG} & 0x02 & 0x04 & 0x06 & --- \\
        \hline
    \end{tabular}

    \caption{Códigos de \textit{NACK} obtidos no modo ocioso (7 casos)}
    \label{qua:ape-i2-ocioso}
    \fonte{Elaborado pelo autor (2026).}
\end{quadro}

% ---------------------------------------------------------------------------- %
\section{Modo de transferência}
\label{sec:ape-i2-transferencia}

\begin{quadro}[htb]
    \centering
    \small
    \begin{tabular}{|p{4.6cm}|c|c|c|p{3.6cm}|}
        \hline
        Comando & Tamanho & CRC & Identificador & Outras condições \\
        \hline
        Comando inexistente (transferência) & --- & --- & --- & 0x20 \\
        \hline
        \texttt{DATA\_PACKET} & 0x21 & 0x23 & 0x25 & Carga útil acima de 256 \textit{bytes}: 0x27 \\
        \hline
        \texttt{STOP\_TRANSFER} & 0x22 & 0x24 & 0x26 & --- \\
        \hline
        \texttt{SET\_DEBOUNCE} & 0x2C & 0x2D & 0x2E & Nível 11 (fora de 1 a 10): 0x2F \\
        \hline
    \end{tabular}

    \caption{Códigos de \textit{NACK} obtidos no modo de transferência (12 casos)}
    \label{qua:ape-i2-transferencia}
    \fonte{Elaborado pelo autor (2026).}
\end{quadro}

% ---------------------------------------------------------------------------- %
\section{Modo de depuração}
\label{sec:ape-i2-depuracao}

\begin{quadro}[htb]
    \centering
    \small
    \begin{tabular}{|p{4.6cm}|c|c|c|p{3.6cm}|}
        \hline
        Comando & Tamanho & CRC & Identificador & Outras condições \\
        \hline
        Comando inexistente (depuração) & --- & --- & --- & 0x40 \\
        \hline
        \texttt{RESUME\_DEBUG} & 0x41 & 0x44 & 0x47 & --- \\
        \hline
        \texttt{STOP\_DEBUG} & 0x42 & 0x45 & 0x48 & --- \\
        \hline
        \texttt{PAUSE\_DEBUG} & 0x43 & 0x46 & 0x49 & --- \\
        \hline
        \texttt{FORCE\_REQ} & 0x4A & 0x4B & 0x4C & --- \\
        \hline
        \texttt{PUBLISH\_VAR} & 0x4D & 0x4E & 0x4F & 17ª variável (lista cheia): 0x53 \\
        \hline
        \texttt{UNPUBLISH\_VAR} & 0x50 & 0x51 & 0x52 & --- \\
        \hline
        \texttt{WRITE\_MEMORY} & 0x54 & 0x55 & 0x56 & --- \\
        \hline
        \texttt{WRITE\_DR} & 0x57 & 0x58 & 0x59 & --- \\
        \hline
    \end{tabular}

    \caption{Códigos de \textit{NACK} obtidos no modo de depuração (26 casos)}
    \label{qua:ape-i2-depuracao}
    \fonte{Elaborado pelo autor (2026).}
\end{quadro}



% ---------------------------------------------------------------------------- %
% Bibliografia
\backmatter \singlespacing   % espaçamento simples
\bibliographystyle{abntex2-alf}% citação bibliográfica alpha alpha-ime
\bibliography{referencias}  % associado ao arquivo: 'referencias.bib'


\end{document}
